%% toc: P15 164 164.6
\section{Carnap 1950: limits, classical laws, and the Johnson--Carnap continuum}
\label{sec:annex-p15-w02b-p15-029}

The definitions of Chapter~\ref{sec:p15-w02b-p15-029} apply. In the
monadic systems, a state of \(\mathfrak L_N\) is determined by the
\(Q\)-property of each \(a_1,\ldots,a_N\), since identity is fixed.
For a sentence \(i\) with constants \(b_1,\ldots,b_m\), the elements
not denoted by any \(b_l\) are called \emph{unnamed} (relative to
\(i\)); the truth of \(i\) in a state does not depend on which
further constants denote them.

\begin{annexitem}{annex:p15-w02b-p15-029-saturation}{thm:p15-w02b-p15-029-saturation}{Counts above the quantifier rank}
\begin{proof}
For tuples \(\bar x\in\dom(\mathcal A)^l\), \(\bar y\in\dom(\mathcal B)^l\)
and \(k\geq0\), write \((\mathcal A,\bar x)\sim_k(\mathcal B,\bar y)\)
when (i) each \(x_p\) has the same \(Q\)-property as \(y_p\);
\(x_p=x_{p'}\) iff \(y_p=y_{p'}\); and \(x_p=b_l^{\mathcal A}\) iff
\(y_p=b_l^{\mathcal B}\); and (ii) for each \(j\), the numbers of
elements of \(Q_j\) that are neither some \(b_l\) nor some entry of
the tuple are equal in the two structures or both at least \(k\).
We show by induction on formulas \(\psi(v_1,\ldots,v_l)\) of quantifier
rank at most \(k\) that \((\mathcal A,\bar x)\sim_k(\mathcal B,\bar y)\)
implies \(\mathcal A\models\psi[\bar x]\iff\mathcal B\models\psi[\bar y]\).

Atomic formulas are \(t\), \(P(\tau)\), or \(\tau=\tau'\), where each
term is a variable or some \(b_l\). The truth of \(P(\tau)\) is fixed
by the \(Q\)-property of the denoted element, and an identity by the
equality pattern; (i) makes these agree. Negation and conjunction pass
directly, since rank does not increase. The universal quantifier is
\(\neg\exists\neg\), so consider \(\exists v\,\psi\) with
\(\psi\) of rank at most \(k-1\), \(k\geq1\). Suppose
\(\mathcal A\models\psi[\bar x,x]\). If \(x\) is some \(x_p\) or
some \(b_l^{\mathcal A}\), let \(y\) be the corresponding element of
\(\mathcal B\); (i) still holds for the extended tuples, and the sets
in (ii) are unchanged, so the counts are equal or both at least
\(k\geq k-1\). Otherwise \(x\) is an unlisted element of some
\(Q_j\). The \(\mathcal A\)-count for \(Q_j\) is then at least one,
so the \(\mathcal B\)-count is equal to it or at least \(k\geq1\);
choose an unlisted \(y\in Q_j\) in \(\mathcal B\). Condition (i)
holds, since \(y\) is new and has the same \(Q\)-property; in (ii)
both counts for \(Q_j\) fall by one, so equal counts stay equal and
counts at least \(k\) become at least \(k-1\), while other counts are
unchanged. In both cases
\((\mathcal A,\bar xx)\sim_{k-1}(\mathcal B,\bar yy)\), and the
induction hypothesis gives \(\mathcal B\models\psi[\bar y,y]\).
The converse direction is symmetric. The lemma is the case \(l=0\),
\(k=r\): the hypothesis is exactly \(\sim_r\) for empty tuples.
\end{proof}
\end{annexitem}

\begin{annexitem}{annex:p15-w02b-p15-029-monadic-limit}{thm:p15-w02b-p15-029-monadic-limit}{Monadic limits, the reduction, and 110A (5)}
\begin{proof}
Let \(n_0\) be the largest index in \(i\) and choose one constant
\(a_*\) not among \(b_1,\ldots,b_m\). For \(N>\max\{n_0,\operatorname{index}(a_*)\}\)
put \(s=m+1\) and \(T=N-s\). For a word \(v\) of \(Q\)-properties
of \(b_1,\ldots,b_m,a_*\), let \(u(v)\) be its restriction to the
\(b\)'s and \(\langle v\rangle\) the corresponding conjunction. By
the fitting property in \cref{thm:p15-carnap-star-construction},
\(m^*_N(\langle v\rangle)=m^*(\langle v\rangle)\) for all such \(N\).

Let \(E_N\) be the set of states in which each \(Q_j\) contains at
least \(r\) of the \(T\) individuals other than
\(b_1,\ldots,b_m,a_*\). In a state of \(E_N\) every \(Q_j\) contains
at least \(r\) unnamed elements (relative to \(i\)), as does
\(D_{r,u}\). By \cref{thm:p15-w02b-p15-029-saturation}, a state
\(z\in E_N\) with \(b\)-profile \(u\) satisfies \(i\) iff
\(u\in S_i\).

We bound \(c^*_N(\neg E_N,\langle v\rangle)\). Apply
\cref{thm:p15-carnap-predictive} to the division
\(\{Q_j,\neg Q_j\}\), with widths \(1\) and \(\kappa-1\), the
evidence \(\langle v\rangle\) about \(s\) individuals with \(s_j\) of
them in \(Q_j\), and the further list of the \(T\) remaining
individuals. For \(0\leq c\leq T\),
\[
 c^*_N(t_j=c,\langle v\rangle)
 =\binom{s_j+c}{c}
  \frac{\binom{T-c+a}{a}}{\binom{T+b}{b}},\qquad
 a=s-s_j+\kappa-2,\ b=s+\kappa-1=a+s_j+1 .
\]
Since \(\kappa\geq2\), \(a\geq0\). Using
\(\binom{T-c+a}a\leq(T+a)^a/a!\) and
\(\binom{T+b}b\geq T^b/b!\),
\[
 c^*_N(t_j=c,\langle v\rangle)\leq
 \binom{s_j+c}{c}\frac{b!}{a!}\frac{(T+a)^a}{T^{a+s_j+1}}
 \longrightarrow0\quad(T\to\infty).
\]
Summing over the \(\kappa\) values of \(j\) and the \(r\) values
\(c<r\) gives \(\varepsilon_N(v)\defeq c^*_N(\neg E_N,\langle v\rangle)\to0\).

Now \(m^*_N(i)=\sum_vm^*_N(\langle v\rangle\land i)\), a finite sum
over \(\kappa^{m+1}\) words. If \(u(v)\in S_i\), then
\(m^*(\langle v\rangle)(1-\varepsilon_N(v))\leq
m^*_N(\langle v\rangle\land i)\leq m^*(\langle v\rangle)\); if
\(u(v)\notin S_i\), then
\(0\leq m^*_N(\langle v\rangle\land i)\leq
m^*(\langle v\rangle)\varepsilon_N(v)\). Hence
\[
 m^*_N(i)\longrightarrow\sum_{v:\,u(v)\in S_i}m^*(\langle v\rangle)
 =\sum_{u\in S_i}m^*(\langle u\rangle)=m^*(i'),
\]
the middle equality by summing over the profile of \(a_*\), and the
last because the \(\langle u\rangle\) are mutually exclusive. The
limit along all \(N\geq n_0\) is the same, since only a final segment
matters. The sentence \(i'\) is nongeneral with constants among the
\(b_l\). Each \(D_{r,u}\) is a finite structure with \(m+r\kappa\)
elements, and truth of \(i\) in it is decided by the recursive
truth definition with quantifiers ranging over a finite set; so
\(S_i\) and \(i'\) are computable, which proves (1).

For (2), \(m=0\): the only word is empty, so \(i'\) is \(t\) or
\(\neg t\), with \(m^*\)-values \(1\) or \(0\). For (3),
\(\forall x\,M(x)\) has rank one and no constants, and
\(D_{1,\varnothing}\) contains an element of a \(Q\)-property outside
\(M\), which exists since \(w<\kappa\); so the sentence is false there
and the limit is \(0\). For (4), if
\(m^*_\infty(e)>0\) then \(m^*_N(e)>0\) for all large \(N\), so
\(c^*_N(h,e)=m^*_N(e\land h)/m^*_N(e)\); both limits exist by (1),
and the quotient of the limits is the limit of the quotients.
\end{proof}
\end{annexitem}

\begin{annexitem}{annex:p15-w02b-p15-029-ltruth-transfer}{thm:p15-w02b-p15-029-ltruth-transfer}{Transfer of L-truth}
\begin{proof}
A sentence is L-true in a system when it holds in every state
description of that system. A state description of
\(\mathfrak L_\infty\) is a structure whose domain is denumerable,
every element denoted by exactly one constant.

(1) Suppose \(i\) fails in a state \(z\) of \(\mathfrak L_N\) with
\(N\geq\max\{n_0,m+\kappa(r-1)+1\}\). There are \(N-m\) unnamed
elements, so by the pigeonhole principle some \(Q_{j_0}\) contains at
least \(r\) of them. Add denumerably many new elements to
\(Q_{j_0}\). In the resulting structure \(\mathcal B\) the unnamed
counts agree with those of \(z\) except in \(Q_{j_0}\), where both are
at least \(r\). The unnamed elements of \(\mathcal B\) form a
denumerable set; biject it with the indices not used by
\(b_1,\ldots,b_m\) and let those constants denote accordingly. This
is a state description of \(\mathfrak L_\infty\), and by
\cref{thm:p15-w02b-p15-029-saturation} it satisfies \(\neg i\). So
\(i\) is not L-true in \(\mathfrak L_\infty\). This is the first
claim in contrapositive form.

Conversely, suppose \(i\) fails in a state description \(\mathcal B\)
of \(\mathfrak L_\infty\), with unnamed counts
\(t_j\in\{0,1,\ldots\}\cup\{\infty\}\). Since there are infinitely
many unnamed elements and finitely many \(Q_j\), some
\(t_{j_1}=\infty\). Let \(N\geq\max\{n_0,m+\kappa r\}\). Choose
\(t'_j=t_j\) if \(t_j<r\), \(t'_j=r\) if \(r\leq t_j<\infty\) and
\(j\ne j_1\), and \(t'_{j_1}=N-m-\sum_{j\ne j_1}t'_j\), which is at
least \(N-m-\kappa r+r\geq r\). The finite structure with the same
named profiles and \(t'_j\) unnamed elements in \(Q_j\) has
\(N\) elements; naming the unnamed ones by the indices in
\(\{1,\ldots,N\}\) not used by the \(b_l\) gives a state of
\(\mathfrak L_N\). Its counts agree with \(\mathcal B\)'s or are both at
least \(r\), so it satisfies \(\neg i\) by the lemma, and \(i\) is not
L-true in \(\mathfrak L_N\).

(2) Suppose \(i\) is L-true in \(\mathfrak L_\infty\) but fails in some
state of \(\mathfrak L_N\) for infinitely many \(N\). Let
\(\Lang\) consist of the primitive predicates and the constants
\(b_1,\ldots,b_m\) of \(i\). Let \(\Gamma\) contain \(\neg i\), the
sentences \(b_l\ne b_{l'}\) for \(l\ne l'\), and for each
\(n\geq1\) a sentence \(\theta_n\) saying that there are at least
\(n\) elements. A finite \(\Gamma_0\subseteq\Gamma\) mentions
finitely many \(\theta_n\), all implied by the one with the largest
\(n\); a failing state of some \(\mathfrak L_N\) with \(N\geq n\),
viewed as an \(\Lang\)-structure, is a model of \(\Gamma_0\). By
\cref{thm:fol-compactness}, \(\Gamma\) has a model
\(\mathcal M\). By \cref{thm:downward-ls} with
\(A=\{b_1^{\mathcal M},\ldots,b_m^{\mathcal M}\}\), some
\(\mathcal N\preccurlyeq\mathcal M\) contains \(A\) and has at most
\(\aleph_0\) elements, \(\Lang\) being finite. As an elementary
substructure, \(\mathcal N\) satisfies every \(\theta_n\), so it is
infinite, hence denumerable; it satisfies \(\neg i\), and the
\(b_l\) denote distinct elements. Naming its remaining elements
bijectively by the unused indices yields a state description of
\(\mathfrak L_\infty\) in which \(\neg i\) holds, a contradiction.
Both cited theorems are stated under the Axiom of Choice; part (1)
uses no choice principle.
\end{proof}
\end{annexitem}

\begin{annexitem}{annex:p15-w02b-p15-029-t81-transfer}{thm:p15-w02b-p15-029-t81-transfer}{T81-2a for general sentences}
\begin{proof}
\(\operatorname{MC}\) in \(\mathfrak L_\infty\) means that \(e,e'\)
are not L-false there and that at least one of the conditions
(\(c_1\)) \(e\supset h\), (\(c_2\)) \(e'\supset\neg h'\), or
(\(c_3\)) both \((e'\land h')\supset(e\land h)\) and
\(e\supset(h\lor e')\), is L-true there. A holding condition involves
at most two sentences; apply
\cref{thm:p15-w02b-p15-029-ltruth-transfer} to each and take the
larger threshold together with the largest index among the four
sentences. Beyond this \(N_0\), the same condition is L-true in
\(\mathfrak L_N\), and \(e,e'\) are not L-false in \(\mathfrak L_N\)
by hypothesis. So \(\operatorname{MC}(h,e;h',e')\) holds in
\(\mathfrak L_N\), and by the finite theorem
\cref{thm:p15-carnap-comparative} (T81-1a)
\(c_N(h,e)\geq c_N(h',e')\) for every regular \(c_N\) and every
\(N\geq N_0\). If both limits exist, a weak inequality between the
terms of two convergent sequences on a final segment persists in the
limit. This is Carnap's printed argument, with the transfer step now
proved. For monadic systems the threshold is explicit by part (1) of
the transfer theorem, and for \(c^*\) both limits exist when
\(m^*_\infty(e),m^*_\infty(e')>0\) by
\cref{thm:p15-w02b-p15-029-monadic-limit}(4).
\end{proof}
\end{annexitem}

\begin{annexitem}{annex:p15-w02b-p15-029-betting}{thm:p15-w02b-p15-029-betting}{Betting quotients and estimates}
\begin{proof}
(1) \(u_1,u_2\geq0\) and \(u_1+u_2>0\) give
\(0\leq u_1/(u_1+u_2)\leq1\). For \(X_2\) the same bet is a bet on
\(\neg h\) with stake \(u_2\) at risk, so its quotient is
\(u_2/(u_1+u_2)=1-q\).

(2) \(X_1\) wins \(rn\) bets, each paying \(u_2\), and loses
\((1-r)n\), each costing \(u_1\). The balance is
\(rnu_2-(1-r)nu_1=n\bigl(r(u_1+u_2)-u_1\bigr)=n(u_1+u_2)(r-q)\).

(3) Let \(P(k)\) be the probability\(_1\) of exactly \(k\) members of
\(K\) having \(M\), \(\sum_kP(k)=1\). The frequency is \(k/n\) and the
gain is \(n(u_1+u_2)(k/n-q)\), both affine in \(k\). An estimate is a
weighted mean with weights summing to one, so it commutes with affine
maps: \(r'=\sum_kP(k)k/n=m'/n\) and
\(g'=n(u_1+u_2)(r'-q)\). Since \(n(u_1+u_2)>0\), \(g'=0\) iff
\(q=r'\), and then the estimated gain of \(X_2\), which is \(-g'\),
is also zero.

(4) The L-true \(j_1\lor\cdots\lor j_n\) has support one. By induction
on \(k\), additivity over L-exclusive pairs gives the support of
\(j_1\lor\cdots\lor j_k\) as the sum of the supports, since
\(j_1\lor\cdots\lor j_{k-1}\) and \(j_k\) are L-exclusive. With equal
supports \(\sigma\), \(n\sigma=1\), so \(\sigma=1/n\), and a disjunction
of \(k\) of them has support \(k/n\).
\end{proof}
\end{annexitem}

\begin{annexitem}{annex:p15-w02b-p15-029-binomial}{thm:p15-w02b-p15-029-binomial}{Binomial law}
\begin{proof}
Write \(x^{\underline j}=x(x-1)\cdots(x-j+1)\). The evidence \(e\) is
invariant under permutations of the population's names, so by
\cref{thm:p15-carnap-regular-laws} all individual distributions of
the fixed sample with the same counts have the same value. There are
\(s!/\prod_is_i!\) of them, and by \cref{thm:p15-carnap-direct} their
disjunction \(h_{st}\) has value \(\prod_i\binom{n_i}{s_i}/\binom ns\).
Dividing,
\[
 c(h,e)=\frac{\prod_in_i^{\underline{s_i}}}{n^{\underline s}},
\]
which is also the probability of a specified ordered outcome of \(s\)
draws without replacement from the population.

(A) For \(x\geq j\geq1\),
\((x-j+1)^j\leq x^{\underline j}\leq x^j\). Hence, using
\(n^s=\prod_in^{s_i}\),
\[
 c(h,e)\geq\frac{\prod_i(n_i-s_i+1)^{s_i}}{n^s}
 =\prod_ir_i^{s_i}\Bigl(1-\frac{s_i-1}{n_i}\Bigr)^{s_i},\qquad
 c(h,e)\leq\frac{\prod_in_i^{s_i}}{(n-s+1)^s}
 =\prod_ir_i^{s_i}\Bigl(1-\frac{s-1}n\Bigr)^{-s}.
\]
For two categories the joint law of the sample's categories is the
law of the first \(s\) ordered draws without replacement; by
\cref{thm:p15-finite-urn-approximation} it is within
\(s(s-1)/(2n)\) in total variation of the independent law with
parameter \(r_1\). The sample count is a function of the outcomes, so
its two laws are no farther apart.

(B) \(c(h,e)=\prod_i\prod_{l<s_i}\frac{n_i-l}n\Big/\prod_{l<s}\frac{n-l}n\).
With \(s\) fixed and \(n_i/n\to r_i\), each factor
\((n_i-l)/n\to r_i\) and \((n-l)/n\to1\), giving
\(\prod_ir_i^{s_i}\). Multiplying by the fixed number
\(s!/\prod_is_i!\) gives the limit for \(h_{st}\), and \(p=2\) gives the
binomial form.

(d) For \(1\leq s_1\leq s\),
\[
 \frac{B(s_1)}{B(s_1-1)}=\frac{(s-s_1+1)r_1}{s_1r_2},
\]
which exceeds, equals, or is less than one according as \(s_1<q\),
\(s_1=q\), or \(s_1>q\), since
\((s-s_1+1)r_1>s_1(1-r_1)\iff(s+1)r_1>s_1\). As
\(0<q<s+1\), the values increase up to \(\lfloor q\rfloor\) and
decrease afterwards when \(q\) is not an integer; when \(q\) is an
integer, \(1\leq q\leq s\) and \(B(q-1)=B(q)\) are the two maxima. If
\(sr_1\) is an integer, \(q=sr_1+r_1\) is not, and
\(\lfloor q\rfloor=sr_1\).
\end{proof}
\end{annexitem}

\begin{annexitem}{annex:p15-w02b-p15-029-bernoulli}{thm:p15-w02b-p15-029-bernoulli}{Bernoulli's theorem}
\begin{proof}
(e) Since \(0<r_1<1\), \(n\geq2\). Let \(S\) be the sample count.
By \cref{thm:p15-carnap-frequency-reliability}, under any symmetric
\(c\), the \(c\)-mean of \(S\) is \(sr_1\) and its estimated square
error is \(J^2=sr_1r_2(n-s)/(n-1)\). Let
\(F=\{x:|x-sr_1|>gs\}\). Then
\[
 c(S\in F,e)=\sum_{x\in F}c(S=x,e)
 \leq\sum_x c(S=x,e)\frac{(x-sr_1)^2}{g^2s^2}
 =\frac{J^2}{g^2s^2}
 =\frac{r_1r_2}{sg^2}\frac{n-s}{n-1}.
\]
The complementary event is \(|s_1/s-r_1|\leq g\). The second inequality
uses \((n-s)/(n-1)\leq1\) and \(r_1r_2\leq1/4\). Given
\(g,\varepsilon\), any \(s'>1/(4g^2\varepsilon)\) serves for every
population size.

(a)--(d) \Cref{thm:p15-stirling-general-binomial} with \(N=s\),
\(p=r_1\) gives the uniform local statement and convergence in
distribution of \(Z_s=(S-sr_1)/\sigma\) to the standard normal law.
Since \(\Phi\) is continuous,
\(P(u_1\leq Z_s\leq u_2)\to\Phi(u_2)-\Phi(u_1)\) for fixed
\(u_1<u_2\). Widening by \(\tfrac12\) changes the endpoints to
\(u_l\mp1/(2\sigma)\); for any \(\eta>0\) these lie eventually in
\([u_1-\eta,u_2+\eta]\), so the widened value is squeezed between the
unwidened value and \(P(u_1-\eta\leq Z_s\leq u_2+\eta)\), whose limit
is \(\Phi(u_2+\eta)-\Phi(u_1-\eta)\); letting \(\eta\to0\) gives the
same limit. Taking \(u_1=-u\), \(u_2=u\) gives \(2\Phi(u)-1\), and
the complement gives \(2(1-\Phi(u))\).

For the exact finite-population values with \(n_1/n=r_1\), the total
variation bound of \cref{thm:p15-w02b-p15-029-binomial}(A) shows that
every interval value differs from its binomial value by at most
\(s(s-1)/(2n)\to0\) when \(s^2/n\to0\). For single values,
\(\sigma\) times the difference is at most
\(\sigma s^2/(2n)\leq s^{5/2}/(4n)\), since \(\sigma\leq\sqrt s/2\);
this tends to zero when \(s^{5/2}/n\to0\).
\end{proof}
\end{annexitem}

\begin{annexitem}{annex:p15-w02b-p15-029-variety}{thm:p15-w02b-p15-029-variety}{Variety under \(c^*\)}
\begin{proof}
The \(Q\)-properties are \(KW,K\neg W,\neg KW,\neg K\neg W\), each of
width one; \(W\) has width two and \(\neg K\) width two. Both
evidences report full profiles of \(s=100\) individuals, all \(W\);
let \(x,y\) be the numbers of \(KW\) and \(\neg KW\) individuals, so
\(x+y=100\).

(1) Given either evidence, \(l\) holds iff all \(T=N-100\) further
individuals are \(W\). Apply \cref{thm:p15-carnap-predictive} with
the \(Q\)-division: the event is the disjunction of the count reports
with \(t_{K\neg W}=t_{\neg K\neg W}=0\) and
\(t_{KW}=\alpha\), \(t_{\neg KW}=T-\alpha\). Hence
\[
 c^*_N(l,e)=\binom{T+103}{T}^{-1}\sum_{\alpha=0}^T
 \binom{x+\alpha}{\alpha}\binom{y+T-\alpha}{T-\alpha}
 =\binom{x+y+T+1}{T}\Big/\binom{T+103}{T},
\]
where the sum is the coefficient of \(z^T\) in
\((1-z)^{-(x+1)}(1-z)^{-(y+1)}=(1-z)^{-(x+y+2)}\). The value depends
only on \(x+y=100\). The singular formula gives
\(c^*(W(a),e)=(100+2)/(100+4)=51/52\) for both.

(2) By the product rule and the singular formula,
\[
 c^*(W(a),e\land\neg K(a))
 =\frac{c^*((\neg K\land W)(a),e)}{c^*(\neg K(a),e)}
 =\frac{(y+1)/104}{(y+2)/104}=\frac{y+1}{y+2},
\]
which is \(1/2\) for \(y=0\) and \(51/52\) for \(y=50\).
\end{proof}
\end{annexitem}

\begin{annexitem}{annex:p15-w02b-p15-029-restricted-law}{thm:p15-w02b-p15-029-restricted-law}{The restricted law}
\begin{proof}
\(l'\) is the report that all \(T=N-s\) individuals outside the sample
are \(\neg M_1\). Apply \cref{thm:p15-carnap-predictive} to the
division \((M_1,\neg M_1)\) with widths \((w_1,\kappa-w_1)\),
sample counts \((s_1,s-s_1)\), and new counts \((0,T)\):
\[
 c^*_N(l',e')=\binom{s_1+w_1-1}{0}
 \binom{s-s_1+T+\kappa-w_1-1}{T}\Big/\binom{s+T+\kappa-1}{T},
\]
which is the display. Put \(y=N-s_1+\kappa-w_1-1\), the upper entry of
the numerator; \(y\geq T\) since \(s-s_1+\kappa-w_1-1\geq0\). For
\(T\geq1\), \(\binom yT\) is strictly increasing in \(y\geq T-1\),
because \(\binom{y+1}T/\binom yT=(y+1)/(y+1-T)>1\) when
\(\binom yT>0\); so the value strictly decreases as \(s_1\) increases.
For the limit write the value as \(\binom{T+a}a/\binom{T+b}b\) with
\(a=s-s_1+\kappa-w_1-1\) and \(b=s+\kappa-1=a+s_1+w_1\); then
\(b>a\), and the bounds used in
\cref{thm:p15-w02b-p15-029-monadic-limit} give a quantity of order
\(T^{a-b}\to0\).

For the second form put \(k=s_1+w_1\); then \(k\leq s+\kappa-1\)
because \(s_1\leq s\) and \(w_1\leq\kappa-1\). The first form equals
\[
 \frac{(N+\kappa-1-k)!}{(N-s)!\,(s+\kappa-1-k)!}\cdot
 \frac{(N-s)!\,(s+\kappa-1)!}{(N+\kappa-1)!}
 =\frac{(s+\kappa-1)!\,(N+\kappa-1-k)!}{(s+\kappa-1-k)!\,(N+\kappa-1)!},
\]
and so does
\(\binom{s+\kappa-1}k/\binom{N+\kappa-1}k\) after the factors
\(k!\) cancel.
\end{proof}
\end{annexitem}

\begin{annexitem}{annex:p15-w02b-p15-029-law-approx}{thm:p15-w02b-p15-029-law-approx}{Formula (11)}
\begin{proof}
By \cref{thm:p15-carnap-universal}, applied to the predicate
\(\neg M_1\) of width \(\kappa-w_1\in(0,\kappa)\) which all \(s\)
reported individuals satisfy,
\[
 c^*_N(l,e)=\frac{(s+\kappa-w_1)^{\overline{N-s}}}{(s+\kappa)^{\overline{N-s}}}
 =\frac{(N+\kappa-w_1-1)!\,(s+\kappa-1)!}{(s+\kappa-w_1-1)!\,(N+\kappa-1)!}
 =\binom{s+\kappa-1}{w_1}\Big/\binom{N+\kappa-1}{w_1},
\]
since \(l\) is L-equivalent to \(\forall x\neg M_1(x)\). Cancelling
the factorials,
\(c^*_N(l,e)=\prod_{j=0}^{w_1-1}(s+a_j)/(N+a_j)\) with
\(a_j=\kappa-1-j\), so \(1\leq\kappa-w_1\leq a_j\leq\kappa-1\).

For \(a\geq0\) and \(N\geq s\), \((s+a)/(N+a)\geq s/N\), because
\(N(s+a)-s(N+a)=a(N-s)\geq0\). Hence
\[
 \rho_N=\prod_{j=0}^{w_1-1}\frac{1+a_j/s}{1+a_j/N}\geq1,
\]
and each factor increases with \(N\), with limit \(1+a_j/s\). So
\(\sup_{N\geq s}\rho_N=\prod_j(1+a_j/s)\). By \(1+x\leq e^x\),
this is at most \(\exp(\sum_ja_j/s)\leq\exp(w_1(\kappa-1)/s)\), which
gives the upper bound. By \(\prod_j(1+x_j)\geq1+\sum_jx_j\) for
\(x_j\geq0\), it is at least
\[
 1+\frac1s\sum_{j=0}^{w_1-1}(\kappa-1-j)
 =1+\frac{w_1(2\kappa-w_1-1)}{2s}\geq1+\frac{w_1(\kappa-1)}{2s},
\]
using \(w_1\leq\kappa-1\). Since \(\rho_N\geq1\),
\(\sup_N\lvert\rho_N-1\rvert\) lies between \(x/2\) and \(e^x-1\) for
\(x=w_1(\kappa-1)/s\), so it tends to zero exactly when \(x\) does.
Finally \(x\leq\kappa^2/s\), and \(x\leq W\kappa/s\) when
\(w_1\leq W\).
\end{proof}
\end{annexitem}

\begin{annexitem}{annex:p15-w02b-p15-029-law-update}{thm:p15-w02b-p15-029-law-update}{A conforming and a violating new instance}
\begin{proof}
Write \(F(\sigma,k)=\binom{\sigma+\kappa-1}k/\binom{N+\kappa-1}k\).
By \cref{thm:p15-w02b-p15-029-restricted-law}, \(c^*_N(l',e')=F(s,k)\).
The evidence \(e'\land\neg M_1(a)\) reports \(s+1\leq N\) individuals
with \(s_1\) negative instances, and \(l''\) is the restricted law for
that sample, so its value is \(F(s+1,k)\); for \(e'\land M_1(a)\) the
negative count is \(s_1+1\), so the value is \(F(s+1,k+1)\). With
\(x=s+\kappa\) and \(y=N+\kappa-1\),
\[
 \frac{F(s+1,k)}{F(s,k)}=\frac{\binom xk}{\binom{x-1}k}=\frac x{x-k},
 \qquad
 \frac{F(s+1,k+1)}{F(s,k)}
 =\frac{\binom x{k+1}}{\binom{x-1}k}\cdot\frac{\binom yk}{\binom y{k+1}}
 =\frac x{k+1}\cdot\frac{k+1}{y-k},
\]
where \(x-k=s-s_1+\kappa-w_1\geq1\) and \(y-k\geq N-s_1>0\). These
are the two displayed ratios. Since \(k\geq w_1\geq1\), the first
exceeds one. As \(l'\) is L-equivalent to \(l''\land\neg M_1(a)\), the
multiplication law of \cref{thm:p15-carnap-regular-laws} gives
\(c^*(l',e')=c^*(\neg M_1(a),e')\,c^*(l'',e'\land\neg M_1(a))\), so
the first ratio is \(1/c^*(\neg M_1(a),e')\). If \(s_1\leq S\), then
\(k/(x-k)\leq(S+\kappa)/(s-S)\to0\). The second ratio is below one iff
\(s+\kappa<N+\kappa-1-k\), that is \(N>s+k+1\), and it tends to zero
as \(N\to\infty\).
\end{proof}
\end{annexitem}

\begin{annexitem}{annex:p15-w02b-p15-029-instance-limits}{thm:p15-w02b-p15-029-instance-limits}{When instance confirmation tends to one}
\begin{proof}
By (16), \(1-c^*_i(l,e)=(s_1+w_1)/(s+\kappa)\), which is at most
\((S+w_1)/(s+\kappa)\) and does not involve \(N\) or \(s_2\). The
fraction tends to zero iff both \(w_1/(s+\kappa)\) and
\(s_1/(s+\kappa)\) do, that is, iff \(s\to\infty\) and
\(s_1/s\to0\). By (17), \(1-c^*_{qi}=x/(x+y)\) with
\(x=s_1+w_1\geq1\) and \(y=s_2+w_2\); this tends to zero iff
\(x/y\to0\). If \(s_1\leq S\), then \(1\leq x\leq S+w_1\), so
\(x/y\to0\) iff \(y\to\infty\) iff \(s_2\to\infty\). If \(s_1,s_2\)
are fixed the value does not change with \(s\), and it is below one
because \(x\geq1\).
\end{proof}
\end{annexitem}

\begin{annexitem}{annex:p15-w02b-p15-029-null-fitting}{thm:p15-w02b-p15-029-null-fitting}{Null confirmation and fitting}
\begin{proof}
(1) \(R_N(t\land j)=R_N(j)\) and \(m_N(t)=1\), so
\(c_N(j,t)=m_N(j)\). The evidence \(t\) is L-false in no system, so
\(c_N(j,t)\) is defined for exactly the \(N\) in which \(j\) occurs,
the same index set as for \(m_N(j)\); the two sequences coincide, and
one limit exists iff the other does.

(2) Every state description \(z'\) of \(\mathfrak L_{N+1}\) has exactly
one state description of \(\mathfrak L_N\) as a subconjunction: the
conjunction of those of its components whose constants all have index at
most \(N\); call it the restriction of \(z'\). A nongeneral sentence
\(j\) of \(\mathfrak L_N\) is a truth-functional compound of atomic
sentences of \(\mathfrak L_N\), and its identity components are fixed
by the rules. So \(j\) holds in \(z'\) iff it holds in the restriction.
Hence \(R_{N+1}(j)\) is the set of \(z'\) whose restriction lies in
\(R_N(j)\), and by condition (b) of
\cref{def:p15-w02b-p15-029-fitting},
\[
 m_{N+1}(j)=\sum_{z\in R_N(j)}\ \sum_{z'\text{ restricting to }z}m_{N+1}(z')
 =\sum_{z\in R_N(j)}m_N(z)=m_N(j).
\]
If \(j\) is L-false both ranges are empty and both values zero.
Induction gives \(m_{N+n}(j)=m_N(j)\), and \(m_\infty(j)\) is the
limit of a constant sequence.

(3) Every state of \(\mathfrak L_N\) is the restriction of some state of
\(\mathfrak L_{N+n}\), and of some state description of
\(\mathfrak L_\infty\) (choose the remaining atomic sentences
arbitrarily); by the argument of (2), a nonempty \(R_N(e)\) gives
nonempty ranges there. By (2) applied to \(e\) and to \(h\land e\),
\(c_{N+n}(h,e)=m_N(h\land e)/m_N(e)=c_N(h,e)\), and the limit of this
constant sequence is \(c_\infty(h,e)\).
\end{proof}
\end{annexitem}

\begin{annexitem}{annex:p15-w02b-p15-029-regular-extremes}{thm:p15-w02b-p15-029-regular-extremes}{Extreme values and L-concepts}
\begin{proof}
Case (i). Every state has positive mass and the masses sum to one.
(a, b) \(m(i)=0\) iff \(R_N(i)=\varnothing\) iff \(i\) is L-false.
(c, d) \(m(i)=1\) iff no state of positive mass lies outside
\(R_N(i)\), iff \(R_N(i)=\Omega_N\), iff \(i\) is L-true. (e, f) A
sentence is factual iff it is neither L-true nor L-false, so this
follows from (a)--(d). (g) \(R_N(i)\subsetneq R_N(j)\), and a state in
\(R_N(j)\setminus R_N(i)\) has positive mass. (h)--(k) By (a)--(d),
whether \(m(i)\) is \(0\) or \(1\) depends only on whether \(i\) is
L-false or L-true, not on \(m\). (l) \(R_N(i)\) is neither empty nor
all of \(\Omega_N\). Give each state of \(R_N(i)\) the mass
\(r/\lvert R_N(i)\rvert\) and each remaining state
\((1-r)/\lvert\Omega_N\setminus R_N(i)\rvert\); all masses are positive
and sum to one, and \(m''(i)=r\).

Case (ii). Let \(i,j\) be nongeneral and occur in \(\mathfrak L_N\).
By \cref{thm:p15-w02b-p15-029-null-fitting}(2), \(m(i)=m_N(i)\),
where \(m_N\) is the \(N\)-th function of the sequence, and likewise for
\(m'\). By the restriction argument in that proof, applied to state
descriptions of \(\mathfrak L_\infty\), \(i\) holds in a state
description of \(\mathfrak L_\infty\) iff it holds in its restriction to
\(\mathfrak L_N\), and every state of \(\mathfrak L_N\) is such a
restriction. So \(i\) is L-false, L-true or factual in
\(\mathfrak L_\infty\) iff it is so in \(\mathfrak L_N\), and \(i\)
L-implies \(j\) in \(\mathfrak L_\infty\) iff the nongeneral sentence
\(i\supset j\) is L-true there, iff it is L-true in
\(\mathfrak L_N\). Each assertion (a)--(k) therefore reduces to case (i)
for \(m_N\) and \(m'_N\).
\end{proof}
\end{annexitem}

\begin{annexitem}{annex:p15-w02b-p15-029-almost}{thm:p15-w02b-p15-029-almost}{Almost L-true sentences and replacement}
\begin{proof}
(1) Since \(m_N(\neg i)=1-m_N(i)\) in every system,
\(m(\neg i)\) exists iff \(m(i)\) does, and \(m(\neg i)=1-m(i)\). So
\(\neg i\) is almost L-true (not L-true, value one) iff \(i\) is not
L-false and \(m(i)=0\). If \(i\) were nongeneral and almost L-true,
case (ii)(c, d) of \cref{thm:p15-w02b-p15-029-regular-extremes} would
make it L-true; so it is general. It is not L-true by definition, and
not L-false, since an L-false sentence has \(m_N(i)=0\) for all \(N\)
and hence \(m(i)=0\); so it is factual. The almost L-false case follows
by negation. Let \(i\) be almost L-true and let \(i\supset j\) be
L-true in \(\mathfrak L_\infty\). By
\cref{thm:p15-w02b-p15-029-ltruth-transfer}, \(i\supset j\) is L-true
in every \(\mathfrak L_N\) of a final segment, so
\(m_N(i)\leq m_N(j)\leq1\) there; as \(m_N(i)\to1\), \(m_N(j)\to1\).
So \(m(j)=1\), and \(j\) is L-true or almost L-true.

(2) The sequence for \(m^*\) is fitting
(\cref{thm:p15-carnap-star-construction}). By
\cref{thm:p15-w02b-p15-029-monadic-limit}(1) the reduction of
\(\exists x\,Q_1(x)\) is \(t\), so its value is one; it is not L-true,
since a state description of \(\mathfrak L_\infty\) may give every
individual \(Q_2\). By item~(3) of that theorem,
\(m^*_\infty(\forall x\,M(x))=0\), so \(m^*_\infty(\neg\forall x\,M(x))=1\);
it is not L-true, since every individual may have a \(Q\)-property in
\(M\) (\(w>0\)).

(3) Since \(j\) and \(j'\) are sentences, in every state description of
every system the truth value of \(i\) is computed from the truth values
of its subsentences, and replacing occurrences of \(j\) by \(j'\) where
the two have the same truth value does not change the result. Suppose
\(m(j)=m(j')=0\) and put \(d=j\lor j'\). In a state where \(d\) is
false, \(j\) and \(j'\) are both false, so \(i\land\neg d\) and
\(i'\land\neg d\) are L-equivalent in every system. Hence, in every
\(\mathfrak L_N\),
\[
 \lvert m_N(i)-m_N(i')\rvert
 \leq m_N(d)\leq m_N(j)+m_N(j'),
\]
because \(m_N(i)\) and \(m_N(i')\) each differ from the common value
\(m_N(i\land\neg d)\) by at most \(m_N(d)\). In a finite system the
right side is zero. In \(\mathfrak L_\infty\) it tends to zero, so if
\(m(i)\) exists, \(m(i')\) exists and is equal. If
\(m(j)=m(j')=1\), apply the same argument with \(d=\neg(j\land j')\),
whose value is at most \(m_N(\neg j)+m_N(\neg j')\to0\). For the last
claim let \(j\) be atomic, so \(m^*(j)=m^*(\neg j)=1/2\), and let
\(i=j\lor j\); replacing the second occurrence of \(j\) by \(\neg j\)
gives \(i'=j\lor\neg j\), with \(m^*(i)=1/2\) and \(m^*(i')=1\).
\end{proof}
\end{annexitem}

\begin{annexitem}{annex:p15-w02b-p15-029-limit-laws}{thm:p15-w02b-p15-029-limit-laws}{Passage of the laws to the infinite system}
\begin{proof}
Write \(A_y\) for the set of \(N\) such that \(y\) and the other
sentences in question occur in \(\mathfrak L_N\) and \(y\) is not
L-false there; by D56-2, \(c_\infty(x,y)\) is the limit of
\(c_N(x,y)\) along \(A_y\).

(1) All \(r\) values are limits along the same set \(A_y\), and the
relation holds at every \(N\in A_y\). Polynomial expressions are
continuous, and limits preserve equalities and non-strict inequalities.

(2) \(A_{e\land i}\subseteq A_e\). Let \(N\in A_e\). If \(e\land i\)
is not L-false in \(\mathfrak L_N\), then
\(c_N(h\land i,e)=c_N(i,e)\,c_N(h,e\land i)\) by
\cref{thm:p15-carnap-regular-laws}. If it is L-false, then
\(c_N(i,e)=c_N(h\land i,e)=0\). Put \(q=c_\infty(i,e)\),
\(x=c_\infty(h,e\land i)\). If \(q>0\), then \(c_N(i,e)>0\), hence
\(N\in A_{e\land i}\), for all \(N\in A_e\) beyond some \(N_1\); so
beyond \(N_1\) the sets \(A_e\) and \(A_{e\land i}\) coincide, the
product law holds there, and its two factors tend to \(q\) and \(x\).
If \(q=0\), then \(0\leq c_N(h\land i,e)\leq c_N(i,e)\to0\) along
\(A_e\), so \(c_\infty(h\land i,e)=0=qx\).
\end{proof}
\end{annexitem}

\begin{annexitem}{annex:p15-w02b-p15-029-regular-c}{thm:p15-w02b-p15-029-regular-c}{Mixtures, chains, and extreme values}
\begin{proof}
In \(\mathfrak L_\infty\), a hypothesis that a sentence is L-false or
that one sentence L-implies another holds, by
\cref{thm:p15-w02b-p15-029-ltruth-transfer}, in every \(\mathfrak L_N\)
of a final segment; there the finite laws apply, and
\cref{thm:p15-w02b-p15-029-limit-laws} carries each step below to the
limit: steps with one evidence by its part (1), and each use of the
multiplication law by its part (2). This is the route of Carnap's
proofs marked (II), which cite T20-11 and the limit laws T40-21.

(1) Let \(e=e_1\lor e_2\) with \(e_1\land e_2\) L-false. Since
\(h\land e\) is L-equivalent to the disjunction of the L-exclusive
sentences \(h\land e_1\), \(h\land e_2\), additivity gives
\(c(h,e)=c(h\land e_1,e)+c(h\land e_2,e)\), and
\(c(e_1,e)+c(e_2,e)=1\). The multiplication law and the L-equivalence
of \(e\land e_p\) with \(e_p\) give
\(c(h\land e_p,e)=c(e_p,e)\,c(h,e_p)\). So \(c(h,e)\) is a convex
combination of \(c(h,e_1)\) and \(c(h,e_2)\). For T59-3d the same
computation with the L-exclusive \(j_1,\ldots,j_n\) gives
\(c(h,j)=\sum_pc(j_p,j)\,c(h,j_p)\) with \(\sum_pc(j_p,j)=1\).

(2) Let \(k_p=h_1\land\cdots\land h_p\). The multiplication law gives
\(c(k_n,e)=c(k_{n-1},e)\,c(h_n,e\land k_{n-1})\), and induction on
\(n\) gives the product. Since \(k_n\) L-implies \(h_n\),
\(c(h_n,e)\geq c(k_n,e)\).

(3) By \cref{thm:p15-w02b-p15-029-null-fitting}(3) and the proof of
\cref{thm:p15-w02b-p15-029-regular-extremes}(ii), case (ii) reduces to
\(\mathfrak L_N\) for nongeneral sentences, so let the system be
\(\mathfrak L_N\). (a, b) \(c(h,e)=1\) iff \(m(e\land\neg h)=0\) iff
\(e\land\neg h\) is L-false (\cref{thm:p15-w02b-p15-029-regular-extremes})
iff \(e\) L-implies \(h\). (c, d) \(c(h,e)=0\) iff \(m(e\land h)=0\)
iff \(e\land h\) is L-false. (e) If \(h\) is not L-false,
\(c(e,h)=0\) iff \(e\land h\) is L-false iff \(c(h,e)=0\). If
\(c(h,e)=1\), \(e\) L-implies \(h\), so \(\neg h\) L-implies
\(\neg e\), and \(c(\neg e,\neg h)=1\), \(c(e,\neg h)=0\) when
\(\neg h\) is not L-false. (f) Put \(j_1=\neg e\), \(j_2=e\land h\),
\(j_3=e\land\neg h\); they are L-exclusive in pairs and L-disjunct,
and \(j_2,j_3\) are not L-false by hypothesis. If \(j_1\) is L-false,
\(j_2\) is factual, and case (l) of
\cref{thm:p15-w02b-p15-029-regular-extremes} gives a regular \(m\) with
\(m(j_2)=r\); then \(m(e)=1\) and \(c(h,e)=r\). Otherwise choose
\(0<g<1\) and spread the masses \(g\), \((1-g)r\), \((1-g)(1-r)\)
equally over the states of \(R(j_1),R(j_2),R(j_3)\); then
\(c(h,e)=(1-g)r/\bigl((1-g)r+(1-g)(1-r)\bigr)=r\), and different
\(g\) give different functions, since \(c(j_1,t)=g\).

Value one without entailment
(\cref{def:p15-w02b-p15-029-two-exceptions}): \(e\) has no individual
constant, and in \(\mathfrak L_N\) (\(N\geq2\)) its range is the set of
states in which exactly two of \(a_1,\ldots,a_N\) are \(\neg M\). For
individuals \(a_k,a_l\), the transposition of the two names maps each
state to an isomorphic one, fixes \(R_N(e)\), and maps
\(R_N(e\land\neg M(a_k))\) onto \(R_N(e\land\neg M(a_l))\); so
symmetry makes \(m_N(e\land\neg M(a_k))\) independent of \(k\). Each
state of \(R_N(e)\) satisfies \(\neg M(a_k)\) for exactly two \(k\), so
\(\sum_km_N(e\land\neg M(a_k))=2m_N(e)\), hence
\(m_N(e\land\neg M(b))=2m_N(e)/N\) and \(c_N(M(b),e)=1-2/N\). A state
description of \(\mathfrak L_\infty\) in which \(b\) and one other
individual are \(\neg M\) and all others \(M\) satisfies \(e\land\neg
M(b)\), so \(e\) does not L-imply \(M(b)\).

(4) (a) \(c(h\land i,e)\leq c(h,e)=0\). (b) By the multiplication law
\(0=c(h\land i,e)=c(i,e)\,c(h,e\land i)\) with \(c(i,e)>0\). (c)
\(c(\neg h,e)=0\), so \(c(\neg h,e\land i)=0\) by (b), and
\(c(h,e\land i)=1\). (d) \(c(i,e)=c(h\land i,e)+c(\neg h\land i,e)\)
and \(c(\neg h\land i,e)\leq c(\neg h,e)=0\). (f)
\(1=c(\neg h_1\lor h_2,e)\leq c(\neg h_1,e)+c(h_2,e)=1-c(h_1,e)+c(h_2,e)\).
(g) \(h_1\equiv h_2\) L-implies \(h_1\supset h_2\) and
\(h_2\supset h_1\), so both have value one, and (f) applies in both
directions.
\end{proof}
\end{annexitem}

\begin{annexitem}{annex:p15-w02b-p15-029-bayes}{thm:p15-w02b-p15-029-bayes}{The division theorem and Bayes's theorem}
\begin{proof}
In \(\mathfrak L_\infty\) each step is carried to the limit as in the
proof of \cref{thm:p15-w02b-p15-029-regular-c}. Condition (C) makes
every displayed denominator positive.

T60-1. (a) is the multiplication law
\(c(h\land i,e)=c(h,e)\,c(i,e\land h)\) divided by \(c(h,e)\). (b)
Both sides equal \(c(h\land i,e)\), by the multiplication law applied
once with \(h\) first and once with \(i\) first. (c) Divide (b) by
\(c(i,e)\). (d) Divide (c) by \(c(h,e)\). (e) In (a), write
\(c(h,e)=c(h\land i,e)+c(h\land\neg i,e)\).

T60-2 and T60-3. Write (c) for \(h_1\) and for \(h_2\) (respectively,
for \(i_1\) and \(i_2\)) and divide; the common factor \(c(i,e)\)
(respectively \(c(h,e)\)) cancels.

T60-5. By (c), \(c(h_p,e\land i)\,c(i,e)=c(h_p,e)\,c(i,e\land h_p)\)
for \(p=1,2\). Adding and using \(c(h_1,e\land i)+c(h_2,e\land i)=1\)
gives \(c(i,e)=\sum_pc(h_p,e)\,c(i,e\land h_p)\); substituting this
into (c) for \(h_1\) gives the formula, and equal priors cancel.

T60-6. By (c),
\(c(h,e\land i)=c(h,e)\,c(i,e\land h)/c(i,e)=c(i\land h,e)/c(i,e)\).
Since \(e\land i\) L-implies \(h_1\lor\cdots\lor h_n\), the sentence
\(e\land i\) is L-equivalent to the disjunction of the
\(e\land i\land h_p\), which are L-exclusive in pairs; so
\(c(i,e)=c(e\land i,e)=\sum_pc(i\land h_p,e)\) by additivity. This is
the first form; the multiplication law
\(c(i\land h_p,e)=c(h_p,e)\,c(i,e\land h_p)\) gives the second, and
equal priors cancel.
\end{proof}
\end{annexitem}

\begin{annexitem}{annex:p15-w02b-p15-029-bayes-laplace}{thm:p15-w02b-p15-029-bayes-laplace}{Bayes's postulate and the rule of succession}
\begin{proof}
Let \(I(k,j)=\int_0^1\theta^k(1-\theta)^j\,d\theta\). Integration by
parts gives \(I(k,j)=\bigl(j/(k+1)\bigr)I(k+1,j-1)\) for \(j\geq1\), and
\(I(k+j,0)=1/(k+j+1)\); so \(I(k,j)=k!\,j!/(k+j+1)!\).

(ii)\(\Rightarrow\)(i): by \cref{thm:p15-binary-definetti},
\(P(S_n=k)=\binom nkI(k,n-k)=1/(n+1)\).
(i)\(\Rightarrow\)(ii): the mixing law \(\lambda\) has moments
\(\int\theta^n\lambda(d\theta)=P(X_1=\cdots=X_n=1)=P(S_n=n)=1/(n+1)
=\int_0^1\theta^n\,d\theta\) for \(n\geq1\), and both are
probabilities. By the uniqueness in \cref{thm:p15-feller-hausdorff},
\(\lambda\) is uniform.

For the rule, exchangeability makes all words of length \(n+1\) with
\(k+1\) successes equally likely, and \(k+1\) of the \(n+1\)
positions can carry the last success, so
\(P(X_{n+1}=1,S_n=k)=P(S_{n+1}=k+1)\frac{k+1}{n+1}
=\frac{k+1}{(n+1)(n+2)}\). Dividing by \(P(S_n=k)=1/(n+1)\) gives
\((k+1)/(n+2)\).
\end{proof}
\end{annexitem}

\begin{annexitem}{annex:p15-w02b-p15-029-johnson}{thm:p15-w02b-p15-029-johnson}{Permutation and combination postulates}
\begin{proof}
Exchangeability gives all sequences with counts \((n_i)\) the same
probability, and there are \(n!/\prod_in_i!\) of them; additivity
gives the first identity. For a sequence \(w\) of positive probability,
\(P(X_{n+1}=c_i\mid w)=P(wc_i)/P(w)\); exchangeability of
\(X_1,\ldots,X_{n+1}\) and of \(X_1,\ldots,X_n\) makes numerator and
denominator functions of the counts of \(wc_i\) and \(w\), hence of
the counts of \(w\). Under the combination postulate,
\(1/\binom{n+t-1}{t-1}=(t-1)!\,n!/(n+t-1)!\), and dividing by
\(n!/\prod_in_i!\) gives \(\prod_in_i!\,(t-1)!/(n+t-1)!\). The
predictive probability is the quotient of this expression at
\((n_1,\ldots,n_i+1,\ldots,n_t)\) and at \((n_1,\ldots,n_t)\), which
is \((n_i+1)/(n+t)\). The number of count vectors is the number of
ways to place \(t-1\) separators among \(n+t-1\) positions,
\(\binom{n+t-1}{t-1}\).
\end{proof}
\end{annexitem}

\begin{annexitem}{annex:p15-w02b-p15-029-johnson-carnap}{thm:p15-w02b-p15-029-johnson-carnap}{The Johnson--Carnap continuum}
\begin{proof}
Write \(p_i=P(X_1=c_i)\) and \(g_i(n)=f_i(0,n)\). By (1), each
sequence \(w\) has \(P(w)>0\) and \(P(wc_i)>0\), so all predictive
probabilities are positive; in particular \(g_i(n)>0\). Every pair
\((a,n)\) with \(0\leq a\leq n\) is the count of \(c_i\) in some
sequence.

\emph{Step 1: common slope.} Let \(n\geq1\). For \(t\geq3\), take
\(i\ne j\) and a third category \(l\). Compare a sequence with counts
\(n_i=a+1\), \(n_j=0\), \(n_l=n-a-1\) and one with \(n_i=a\),
\(n_j=1\), \(n_l=n-a-1\) (\(0\leq a\leq n-1\)). In each, the
predictive probabilities sum to one, and only the terms for \(i\) and
\(j\) differ; so
\(f_i(a+1,n)-f_i(a,n)=f_j(1,n)-f_j(0,n)\). The left side is thus
independent of \(a\), and equals the right side, which by the same
identity with \(i,j\) exchanged is \(f_i(1,n)-f_i(0,n)\) for every
\(j\ne i\). Hence there is \(d_n\) with
\(f_i(a,n)=g_i(n)+ad_n\) for all \(i\) and \(a\). For \(t=2\) the
affine form is assumed, and
\(f_1(a,n)+f_2(n-a,n)=1\) for \(0\leq a\leq n\) forces the two slopes
to agree. In both cases, summing over a sequence with all \(n\)
observations in one category,
\[
 \textstyle\sum_ig_i(n)+nd_n=1.\tag{\(\ast\)}
\]

\emph{Step 2: the exchange identity.} For a sequence \(w\) of length
\(n\geq0\) with counts \((n_l)\) and \(i\ne j\),
\(P(wc_ic_j)=P(wc_jc_i)\); dividing by \(P(w)>0\),
\[
 f_i(n_i,n)f_j(n_j,n+1)=f_j(n_j,n)f_i(n_i,n+1).
\]
At \(n=0\) this is \(p_ig_j(1)=p_jg_i(1)\), so
\(p_i=g_i(1)/\sum_lg_l(1)\).

\emph{Step 3 for \(t\geq3\), \(n\geq1\).} Write \(g_i=g_i(n)\),
\(g_i'=g_i(n+1)\), \(d=d_n\), \(d'=d_{n+1}\). Using Step 1 at levels
\(n\) and \(n+1\), the identity reads
\((g_i+n_id)(g_j'+n_jd')=(g_j+n_jd)(g_i'+n_id')\). Using a third
category to hold the remaining observations, it holds for
\((n_i,n_j)=(0,0),(1,0),(0,1)\). These give
\[
 g_ig_j'=g_jg_i',\qquad g_id'=dg_i'\quad\text{for all }i.
\]

\emph{Step 3 for \(t=2\), \(n\geq1\).} With \(n_1=a\), \(n_2=n-a\),
the identity is a polynomial equation in \(a\) whose \(a^2\)
coefficients \(-dd'\) cancel; it holds at \(a=0,1\), so the linear and
constant coefficients agree:
\[
 d(g_1'+g_2')=d'(g_1+g_2),\qquad
 g_1(g_2'+nd')=g_1'(g_2+nd).
\]
By \((\ast)\), the first becomes \(d(1-(n+1)d')=d'(1-nd)\), that is
\(d-d'=dd'\).

\emph{Step 4: the degenerate case.} In both cases \(d_n=0\) iff
\(d_{n+1}=0\) (for \(t\geq3\) since \(g_i>0\); for \(t=2\) from
\(d-d'=dd'\)). Suppose \(d_1=0\). Then all \(d_n=0\), and by
\((\ast)\) \(\sum_ig_i(n)=1\) for each \(n\geq1\). For \(t\geq3\),
\(g_ig_j'=g_jg_i'\) makes \(g_i'/g_i\) independent of \(i\), and
summing shows the common ratio is one. For \(t=2\),
\(g_1g_2'=g_1'g_2\) with \(g_1+g_2=g_1'+g_2'=1\) gives
\(g_1=g_1'\). So \(f_i(a,n)=g_i(1)\) for all \(n\geq1\), and by Step
2 \(p_i=g_i(1)\). Every predictive probability equals \(p_i>0\), so
each sequence has probability \(\prod p_{x_l}\): the observations are
independent with common law \((p_i)\).

\emph{Step 5: the continuum.} Suppose \(d_1\ne0\); then every
\(d_n\ne0\). Put \(\alpha_i(n)=g_i(n)/d_n\). For \(t\geq3\),
\(g_id'=dg_i'\) gives \(\alpha_i(n+1)=\alpha_i(n)\). For \(t=2\),
dividing \(d-d'=dd'\) by \(dd'\) gives \(1/d_{n+1}=1/d_n+1\); with
\((\ast)\), \(\alpha_1(n)+\alpha_2(n)=1/d_n-n\) is the same at
\(n\) and \(n+1\); dividing the constant-coefficient identity by
\(dd'\) gives
\(\alpha_1(\alpha_2'+n)=\alpha_1'(\alpha_2+n)\), and substituting
\(\alpha_2=\beta-\alpha_1\), \(\alpha_2'=\beta-\alpha_1'\), with
\(\beta\) that common sum, gives \((\beta+n)(\alpha_1-\alpha_1')=0\).
As \(\beta+n=1/d_n\ne0\), \(\alpha_1'=\alpha_1\), and so also
\(\alpha_2'=\alpha_2\). In both cases \(\alpha_i\defeq\alpha_i(n)\)
is independent of \(n\geq1\). Put \(\alpha=\sum_i\alpha_i\); by
\((\ast)\), \(d_n(\alpha+n)=1\), so
\[
 f_i(a,n)=d_n(a+\alpha_i)=\frac{a+\alpha_i}{n+\alpha}\qquad(n\geq1).
\]
Signs: \(g_i(n)>0\) and \(\alpha_i=g_i(n)/d_n\) does not depend on
\(n\), so every \(d_n\) has the sign of \(\alpha_1\), and all
\(\alpha_i\) share that sign. If that sign were negative, then \(\alpha<0\) and
\(d_n=1/(n+\alpha)>0\) for \(n>-\alpha\), a contradiction. So
\(d_n>0\) and \(\alpha_i>0\). Finally, Step 2 gives
\(p_i=\alpha_id_1/\sum_l\alpha_ld_1=\alpha_i/\alpha\), the formula at
\(n=0\).

Infinite extendability was used only to choose \(n>-\alpha\) in the
sign argument.
\end{proof}
\end{annexitem}

\begin{annexitem}{annex:p15-w02b-p15-029-laplace-family}{thm:p15-w02b-p15-029-laplace-family}{The family of sec.~47A}
\begin{proof}
Since \(s+2m>0\),
\[
 c_m-r=\frac{s(s_1+m)-s_1(s+2m)}{s(s+2m)}=\frac{m(s-2s_1)}{s(s+2m)}
 =\frac{m(1-2r)}{s+2m},
\]
and
\[
 c_m-\tfrac12=\frac{2(s_1+m)-(s+2m)}{2(s+2m)}=\frac{2s_1-s}{2(s+2m)}
 =\frac{s(2r-1)}{2(s+2m)}.
\]
For \(m=0\) the first difference is zero. Let \(m>0\). If
\(r=\tfrac12\) the second difference is zero. Otherwise both
differences are nonzero, and their signs are those of \(1-2r\) and
\(2r-1\), which are opposite; so \(c_m\) lies strictly between \(r\)
and \(\tfrac12\). As \(0\leq r\leq1\), this interval lies in
\([0,1]\) and \(c_m\) is strictly inside it, so \(0<c_m<1\); for
\(r=\tfrac12\), \(c_m=\tfrac12\). If \(r>\tfrac12\), then
\(c_m-r=m(1-2r)/(s+2m)<0\).
\end{proof}
\end{annexitem}

\begin{annexitem}{annex:p15-w02b-p15-029-decision}{thm:p15-w02b-p15-029-decision}{The decision rules and the consequences of Bernoulli's law}
We first prove the weighted inequality of arithmetic and geometric
means in the form used below.

\emph{Lemma.} Let \(x_1,\ldots,x_n>0\) and \(w_1,\ldots,w_n\geq0\)
with \(\sum_kw_k=1\). Then \(\prod_kx_k^{w_k}\leq\sum_kw_kx_k\), with
equality if and only if all \(x_k\) with \(w_k>0\) are equal.

\emph{Proof of the lemma.} Let \(\varphi(x)=x-1-\log x\) on
\((0,\infty)\). Then \(\varphi'(x)=1-1/x\) is negative on \((0,1)\)
and positive on \((1,\infty)\), so by the mean value theorem
\(\varphi\) is strictly decreasing on \((0,1]\) and strictly
increasing on \([1,\infty)\); since \(\varphi(1)=0\),
\(\log x\leq x-1\) with equality only at \(x=1\). Put
\(A=\sum_kw_kx_k>0\). Then
\[
 \sum_kw_k\log\frac{x_k}A\leq\sum_kw_k\Bigl(\frac{x_k}A-1\Bigr)=0,
\]
that is, \(\log\prod_kx_k^{w_k}\leq\log A\), with equality exactly
when \(x_k=A\) for every \(k\) with \(w_k>0\).

\begin{proof}
(1) Let the values form a finite set \(V\) with probabilities
\(q(u)\), and suppose \(u\in V\) implies \(2u_0-u\in V\) with
\(q(2u_0-u)=q(u)\). The map \(u\mapsto2u_0-u\) is a bijection of
\(V\), so
\(\sum_{u}q(u)(u-u_0)=\sum_uq(2u_0-u)(u_0-u)=-\sum_uq(u)(u-u_0)\).
Hence \(\sum_uq(u)(u-u_0)=0\), and since \(\sum_uq(u)=1\) the
estimate \(\sum_uq(u)u\) equals \(u_0\).

(2) Since \(g'_i=f'_i-f_0\), \(g'_1\geq g'_2\) if and only if
\(f'_1\geq f'_2\). With two actions, \(R_4\) chooses \(j_1\) exactly
when \(g'_1\geq g'_2\). If \(j_1\) is favorable, \(j_1\) is the only
maximizer; if unfavorable, \(j_2\) is; if neutral, both are, and \(R_4\)
permits either. This is \(R_4^*\). If \(\bar u(f)=\gamma f+\delta\) with
\(\gamma>0\) on an interval containing \(f_0\) and all \(f_{ik}\), then
\(\bar g'_i=\sum_kq_{ik}\gamma(f_{ik}-f_0)=\gamma g'_i\); a positive
multiple has the same maximizers.

(3) The two possible fortunes are \(f_0+u\) and \(f_0-u\), each with
probability \(\tfrac12\). The estimated gain is
\(\tfrac12u+\tfrac12(-u)=0\). The estimate of utility is
\[
 \bar g'=\tfrac12\bigl[\bar u(f_0+u)-\bar u(f_0)\bigr]
         -\tfrac12\bigl[\bar u(f_0)-\bar u(f_0-u)\bigr].
\]
Apply the hypothesis with the gain \(g=u>0\) at the fortunes
\(f_0-u<f_0\): \(\bar u(f_0+u)-\bar u(f_0)<\bar u(f_0)-\bar u(f_0-u)\).
Hence \(\bar g'<0\).

(4) The function \(f\mapsto\bar u(f)-k\log f\) has derivative zero on
the interval \((0,\infty)\), so by the mean value theorem it is a
constant \(C\). Hence \(\bar u(f)=k\log f+C\) and
\(\bar u(f_1)-\bar u(f_0)=k\log(f_1/f_0)\). As \(k>0\) and \(\log\)
is injective, \(k\log(f_1/f_0)=k\log(f_3/f_2)\) if and only if
\(f_1/f_0=f_3/f_2\).

(5) By (4),
\(\bar g'=\sum_kq_kk\log\bigl((f_0+g_k)/f_0\bigr)
 =k\log\bigl(\prod_k(f_0+g_k)^{q_k}/f_0\bigr)\).
The function \(\bar u=k\log f+C\) is continuous and strictly increasing
from \((0,\infty)\) onto \(\mathbb R\), so the money equivalent exists
and is unique: \(k\log\bigl((f_0+{}^*g')/f_0\bigr)=\bar g'\), that is,
\(f_0+{}^*g'=f_0e^{\bar g'/k}\). Comparing with the expression for
\(\bar g'\) gives \(f_0+{}^*g'=\prod_k(f_0+g_k)^{q_k}\). The map
\(\bar g'\mapsto f_0(e^{\bar g'/k}-1)\) is strictly increasing, so
maximizing \(\bar g'\) over actions and maximizing \({}^*g'\) select the
same actions. By the lemma with \(x_k=f_0+g_k\) and \(w_k=q_k\),
\(f_0+{}^*g'\leq\sum_kq_k(f_0+g_k)=f_0+g'\), with equality exactly when
the \(g_k\) with \(q_k>0\) coincide.

(6) Item (5) with gains \(u\) and \(-u\), each of probability
\(\tfrac12\), gives \({}^*f'=(f_0+u)^{1/2}(f_0-u)^{1/2}\). The two
fortunes differ because \(u>0\), so the lemma gives strict inequality
with their arithmetic mean \(f_0\).

(7) Insuring yields the fortune \(f_0-r>0\) whether or not the house
burns, with utility \(k\log((f_0-r)/f_0)\). Not insuring yields
\(f_0-h\) with probability \(p\) and \(f_0\) with probability \(1-p\),
with estimate of utility \(pk\log((f_0-h)/f_0)\). Since \(k>0\) and
\(\log\) is strictly increasing, insuring has the larger estimate if and
only if
\(\log(f_0-r)>(1-p)\log f_0+p\log(f_0-h)\), that is,
\(f_0-r>f_0^{1-p}(f_0-h)^p\), that is, \(r<r^*\). Further
\(f_0^{1-p}(f_0-h)^p=f_0\bigl((f_0-h)/f_0\bigr)^p=f_0/q\), so
\(r^*=f_0(1-1/q)\).

By the lemma with the unequal values \(f_0\), \(f_0-h\) and weights
\(1-p\), \(p\) in \((0,1)\),
\(f_0^{1-p}(f_0-h)^p<(1-p)f_0+p(f_0-h)=f_0-ph\), so \(r^*>ph\).

For monotonicity, write \(R(f)=f-f^{1-p}(f-h)^p\) on \((h,\infty)\)
and \(y=(f-h)/f\in(0,1)\). Then
\[
 R'(f)=1-(1-p)f^{-p}(f-h)^p-pf^{1-p}(f-h)^{p-1}=1-\bigl[(1-p)y^p+py^{p-1}\bigr].
\]
Apply the lemma to \(a=y^p\) and \(b=y^{p-1}\) with weights \(1-p\)
and \(p\): \(a\ne b\) because \(a/b=y\ne1\), so
\((1-p)a+pb>a^{1-p}b^p=y^{p(1-p)+p(p-1)}=1\). Hence \(R'(f)<0\) and
\(R\) is strictly decreasing.

For the limit, put \(x=h/f\in(0,1)\); then
\(R(f)=h\cdot\bigl(1-(1-x)^p\bigr)/x\). As \(f\to\infty\), \(x\to0\),
and \(\bigl(1-(1-x)^p\bigr)/x\to p\), the derivative of
\(x\mapsto-(1-x)^p\) at \(0\). Hence \(r^*\to ph\).
\end{proof}
\end{annexitem}

\begin{annexitem}{annex:p15-w02b-p15-029-t53}{thm:p15-w02b-p15-029-t53}{T53-1, and the conventions characterize the regular functions}
In \(\mathfrak L_N\), each sentence \(j\) has a range \(R(j)\), the set
of state descriptions in which it holds; \(j\) is L-true if
\(R(j)=\Omega_N\), L-false if \(R(j)=\varnothing\); \(e\) L-implies
\(h\) if \(R(e)\subseteq R(h)\); and L-equivalence is equality of
ranges (\cref{def:p15-carnap-state-range}). A state description \(z\)
has range \(\{z\}\), and the range of a disjunction or conjunction is
the union or intersection of the ranges.
\begin{proof}
(1a) Apply C53-1c with \(t\), \(e\), \(h\) in place of \(e\), \(h\),
\(j\); all terms are defined because \(t\) and \(t\land e\) are not
L-false. This gives \(c(e\land h,t)=c(e,t)\times c(h,t\land e)\). Since
\(t\land e\) and \(e\) have the same range, C53-1a gives
\(c(h,t\land e)=c(h,e)\).

(1b) Divide (1a) by \(c(e,t)\ne0\).

(1c) Induction on \(n\); \(n=1\) is trivial. For \(n\geq2\), put
\(H=h_1\lor\cdots\lor h_{n-1}\). The range of \(e\land H\land h_n\) is
the union of the ranges of \(e\land h_p\land h_n\) for \(p<n\), all
empty; so \(e\land H\land h_n\) is L-false, and C53-1d gives
\(c(H\lor h_n,e)=c(H,e)+c(h_n,e)\). The hypothesis of (1c) holds for
\(h_1,\ldots,h_{n-1}\), so by induction \(c(H,e)=\sum_{p<n}c(h_p,e)\).
Any bracketing of the disjunction has the same range, so C53-1b
removes the dependence on bracketing.

(1d) \(h\lor\neg h\) has range \(\Omega_N\), so by C53-1b and C53-2
\(c(h\lor\neg h,e)=c(t,e)=1\); and \(e\land h\land\neg h\) is
L-false, so C53-1d gives \(c(h\lor\neg h,e)=c(h,e)+c(\neg h,e)\).

(1e) By (1d) with \(h=t\) and C53-2, \(c(\neg t,e)=1-1=0\).

(1f) An L-false \(h\) has the range of \(\neg t\), so C53-1b and (1e)
give \(c(h,e)=0\).

(2) Suppose first that \(c\) is regular, based on \(m\) with
\(m(z)>0\) and \(\sum_zm(z)=1\). For \(e\) not L-false,
\(m(e)=\sum_{z\in R(e)}m(z)>0\). C53-1a and b hold because
\(m(e\land h)\) and \(m(e)\) depend only on ranges. For C53-1c,
whenever \(e\) and \(e\land h\) are not L-false,
\[
 c(h,e)c(j,e\land h)=\frac{m(e\land h)}{m(e)}\cdot
 \frac{m(e\land h\land j)}{m(e\land h)}=\frac{m(e\land h\land j)}{m(e)}
 =c(h\land j,e).
\]
For C53-1d, if \(e\land h\land j\) is L-false then
\(R(e\land h)\) and \(R(e\land j)\) are disjoint with union
\(R(e\land(h\lor j))\), and \(m\) is additive over disjoint sets of
states. C53-2 holds as \(R(e\land t)=R(e)\), and C53-3 as
\(c(z,t)=m(z)/m(t)=m(z)>0\).

Conversely, suppose \(c\) satisfies the conventions and is defined
exactly when the evidence is not L-false. Put \(m(z)=c(z,t)\) for each
state description \(z\); \(m(z)>0\) by C53-3. Distinct state
descriptions \(z,z'\) have disjoint ranges, so \(t\land z\land z'\) is
L-false; the disjunction of all state descriptions has range
\(\Omega_N\). By (1c), C53-1b and C53-2,
\(\sum_zm(z)=c(t,t)=1\). So \(m\) is a regular \(m\)-function. For a
sentence \(j\) that is not L-false, \(j\) and the disjunction of the
state descriptions in \(R(j)\) have the same range, so by C53-1b and
(1c) \(c(j,t)=\sum_{z\in R(j)}m(z)=m(j)\); for L-false \(j\),
\(c(j,t)=0=m(j)\) by (1f). For \(e\) not L-false, \(c(e,t)=m(e)>0\),
and (1b) gives \(c(h,e)=m(e\land h)/m(e)\). Hence \(c\) is the regular
\(c\)-function based on \(m\). If \(c\) is also based on a regular
\(m'\), then \(c(z,t)=m'(z)/m'(t)=m'(z)\), so \(m'=m\).
\end{proof}
\end{annexitem}

\begin{annexitem}{annex:p15-w02b-p15-029-axiom-systems}{thm:p15-w02b-p15-029-axiom-systems}{What the comparison establishes}
The range notation of the preceding proof is used.
\begin{proof}
(1) Let \(c\) be based on the regular \(m\), and let \(e\) be not
L-false, so \(m(e)>0\). (a) \(R(e\land h)\subseteq R(e)\) and \(m\geq0\)
give \(0\leq m(e\land h)\leq m(e)\). (b) If \(R(e)\subseteq R(h)\), then
\(R(e\land h)=R(e)\) and \(c(h,e)=1\); (c) is the case
\(R(h)=\Omega_N\). (e) If \(R(e\land h)=\varnothing\), \(c(h,e)=0\).
(h), (i) The values depend only on ranges. (k) \(R(e\land(h\lor i))\) is
the union of \(R(e\land h)\) and \(R(e\land i)\), whose intersection is
\(R(e\land h\land i)\); finite additivity of \(m\) over sets of states
gives
\(m(e\land(h\lor i))=m(e\land h)+m(e\land i)-m(e\land h\land i)\);
divide by \(m(e)\). (l) is (k) with \(c(h\land i,e)=0\). (n) is the
computation for C53-1c in the proof of
\cref{thm:p15-w02b-p15-029-t53}(2), under condition (C) that
\(e\land h\) is not L-false, together with (i) for the second form.
(p) Apply (l) to \(h\) and \(\neg h\) (with \(c(h\land\neg h,e)=0\) by
(e)) and use (c) for \(h\lor\neg h\). T59-2a: \(R(e\land i)\) is the
disjoint union of \(R(e\land h\land i)\) and \(R(e\land\neg h\land i)\).
T59-2h: if \(R(e\land h)\subseteq R(j)\), then
\(R(e\land h\land j)=R(e\land h)\).

(2) Let \(q\) be any sentence and \(p=q\land\neg q\). Then \(p\)
L-implies \(q\), since \(R(p)=\varnothing\subseteq R(q)\), and likewise
\(p\) L-implies \(\neg q\). The first condition gives
\(P(q\mid p)=1\), the second \(P(q\mid p)=0\), and the pair
\((q,p)\) has one value; so \(1=0\). If L-false evidence is excluded,
both conditions apply to a pair only when \(R(e)\subseteq R(q)\cap
R(\neg q)=\varnothing\), which is excluded; and any regular \(c\)
satisfies both conditions by (1b) and (1e), so the restricted conditions
are jointly satisfiable.

(3) By \cref{thm:p15-w02b-p15-029-regular-c}(3), items (a)--(d) of
T59-5, \(c(h,e)=1\) only if \(e\) L-implies \(h\) and \(c(h,e)=0\)
only if \(e\land h\) is L-false, that is, \(e\) L-implies \(\neg h\);
neither holds, so \(0<c(h,e)<1\) for every regular \(c\). By the same
theorem, item (f), each \(r\in(0,1)\) equals \(c(h,e)\) for some
regular \(c\). That regular \(c\) satisfies every statement valid for
all regular \(c\)-functions, so no such statement excludes the value
\(r\).
\end{proof}
\end{annexitem}

\begin{annexitem}{annex:p15-w02b-p15-029-kemeny}{thm:p15-w02b-p15-029-kemeny}{Coherence and the five axioms}
Write \(\mathbf1_j(z)\) for \(1\) if \(z\in R(j)\) and \(0\)
otherwise, and \(E=e_1\lor\cdots\lor e_n\) for a book
\((h_k,e_k,S_k)_{k\leq n}\); \(E\) is not L-false. The net gain is
\(G(z)=\sum_kS_k\mathbf1_{e_k}(z)\bigl(\mathbf1_{h_k}(z)-c(h_k,e_k)\bigr)\),
and \(G=0\) outside \(R(E)\).
\begin{proof}
\emph{Coherence implies the axioms.} In each case we exhibit a book
whose gain is negative on all of \(R(E)\) when the axiom fails.
(1) If \(c(h,e)<0\), the single bet with \(S=-1\) has
\(G=c(h,e)-\mathbf1_h\leq c(h,e)<0\) on \(R(e)\); if \(c(h,e)>1\), the
bet with \(S=1\) has \(G=\mathbf1_h-c(h,e)\leq1-c(h,e)<0\) on
\(R(e)\). (3) If \(e\) L-implies \(h\), then \(\mathbf1_h=1\) on
\(R(e)\); if \(c(h,e)<1\), the bet with \(S=-1\) has \(G=c(h,e)-1<0\)
there. (2) If \(R(h)=R(h')\) and \(R(e)=R(e')\), the book with stake
\(1\) on \((h,e)\) and \(-1\) on \((h',e')\) has
\(G=c(h',e')-c(h,e)\) on \(R(e)=R(E)\); if the values differ, this book
or its negative has constant negative gain. (4) If \(e\) L-implies
\(\neg(h\land h')\), then \(\mathbf1_{h\lor h'}=\mathbf1_h+\mathbf1_{h'}\)
on \(R(e)\); the book with stake \(1\) on \((h\lor h',e)\) and \(-1\)
on \((h,e)\) and on \((h',e)\) has
\(G=c(h,e)+c(h',e)-c(h\lor h',e)\) on \(R(e)\). (5) Let \(h\land e\)
be not L-false and put \(a=c(h,e)\), \(b=c(h',h\land e)\),
\(x=c(h\land h',e)\). Take stake \(1\) on \((h\land h',e)\), \(-b\) on
\((h,e)\) and \(-1\) on \((h',h\land e)\). Here \(E\) has the range of
\(e\), and on \(R(e)\), \(\mathbf1_{h\land e}=\mathbf1_h\), so
\[
 G=(\mathbf1_h\mathbf1_{h'}-x)-b(\mathbf1_h-a)-\mathbf1_h(\mathbf1_{h'}-b)=ab-x.
\]
In (2), (4) and (5) the gain is a constant on \(R(E)\), so a nonzero
value gives a sure loss for the book or its negative.

\emph{The axioms imply coherence.} Assume (1)--(5) and fix a book.
First, if \(E\) L-implies \(\neg j\), then \(c(j,E)=0\): \(j\lor j\)
and \(j\) have the same range, and \(E\) L-implies
\(\neg(j\land j)\), so (2) and (4) give \(c(j,E)=2c(j,E)\). For a
state description \(z\in R(E)\) put \(P(z)=c(z,E)\geq0\) (axiom (1)).
For any sentence \(j\), \(j\) has the range of
\((j\land E)\lor(j\land\neg E)\), and \(E\) L-implies
\(\neg\bigl((j\land E)\land(j\land\neg E)\bigr)\), so (2) and (4) give
\(c(j,E)=c(j\land E,E)+c(j\land\neg E,E)\), and the second term is
\(0\). If \(R(j\land E)\) is empty, \(c(j\land E,E)=0\) as well;
otherwise \(j\land E\) has the range of the disjunction of the states
in \(R(j\land E)\), which \(E\) makes exclusive in pairs, and (2) with
(4), applied repeatedly, gives \(c(j\land E,E)=\sum_{z\in R(j\land E)}P(z)\).
Write \(P(j)\) for this sum. By (3), \(P(E)=c(E,E)=1\).

For each \(k\), \(e_k\land E\) has the range of \(e_k\), which is not
empty. Axiom (5) with \(E\), \(e_k\), \(h_k\) in place of \(e\),
\(h\), \(h'\), and then (2), give
\(c(e_k\land h_k,E)=c(e_k,E)\,c(h_k,e_k)\), that is,
\(P(e_k\land h_k)=P(e_k)\,c(h_k,e_k)\). Hence
\[
 \sum_{z\in R(E)}P(z)G(z)=\sum_kS_k\bigl(P(e_k\land h_k)-c(h_k,e_k)P(e_k)\bigr)=0.
\]
If \(G(z)<0\) for every \(z\in R(E)\), then, because the \(P(z)\) are
nonnegative with sum \(1\), the left side would be negative. So no book
yields a sure loss.
\end{proof}
\end{annexitem}

\begin{annexitem}{annex:p15-w02b-p15-029-dirichlet}{thm:p15-w02b-p15-029-dirichlet}{The Dirichlet identification}
For \(x>0\) and an integer \(r\geq0\) write
\(x^{(r)}=x(x+1)\cdots(x+r-1)\), with \(x^{(0)}=1\); then
\(\Gamma(x+r)=x^{(r)}\Gamma(x)\) by \(\Gamma(y+1)=y\Gamma(y)\).
\begin{proof}
(1) By the chain rule for conditional probabilities,
\(P(X_1=e_1,\ldots,X_n=e_n)\) is the product over \(l=0,\ldots,n-1\)
of \(P(X_{l+1}=e_{l+1}\mid X_1=e_1,\ldots,X_l=e_l)\), each conditioning
event having positive probability. By hypothesis the factor for step
\(l\) is \((m+\alpha_i)/(l+\alpha)\) when \(e_{l+1}=c_i\) and \(c_i\)
has occurred \(m\) times among \(e_1,\ldots,e_l\). The denominators
run through \(\alpha,\alpha+1,\ldots,\alpha+n-1\). For each \(i\), the
successive occurrences of \(c_i\) contribute the numerators
\(\alpha_i,\alpha_i+1,\ldots,\alpha_i+n_i-1\). Hence the product is
\(\prod_i\alpha_i^{(n_i)}/\alpha^{(n)}\).

(2) \emph{The Dirichlet integral.} Let \(\beta_1,\ldots,\beta_t>0\),
\(\beta=\sum_i\beta_i\), and let \(\Delta\) be the set of
\((p_1,\ldots,p_{t-1})\) with \(p_i>0\) and \(\sum_{i<t}p_i<1\), with
\(p_t=1-\sum_{i<t}p_i\). We show
\[
 D(\beta_1,\ldots,\beta_t)\defeq\int_\Delta\prod_{i=1}^tp_i^{\beta_i-1}\,dp_1\cdots dp_{t-1}
 =\frac{\prod_i\Gamma(\beta_i)}{\Gamma(\beta)}.
\]
By Tonelli's theorem,
\(\prod_i\Gamma(\beta_i)=\int_{(0,\infty)^t}\prod_ix_i^{\beta_i-1}e^{-\sum_ix_i}\,dx\).
The map \((s,p_1,\ldots,p_{t-1})\mapsto x\) with \(x_i=sp_i\) for
\(i<t\) and \(x_t=sp_t\) is a bijection from \((0,\infty)\times\Delta\)
onto \((0,\infty)^t\), with inverse \(s=\sum_ix_i\), \(p_i=x_i/s\). Its
Jacobian matrix has rows \((p_i,s\,\delta_{i1},\ldots,s\,\delta_{i,t-1})\)
for \(i<t\) and \((p_t,-s,\ldots,-s)\). Adding the first \(t-1\) rows to
the last turns it into \((1,0,\ldots,0)\), since \(\sum_ip_i=1\);
expanding along it leaves \(\pm\det(sI_{t-1})\), so the absolute
Jacobian is \(s^{t-1}\). By the change of variables formula and Tonelli,
\[
 \prod_i\Gamma(\beta_i)=\int_0^\infty s^{\beta-t}e^{-s}s^{t-1}\,ds\int_\Delta\prod_ip_i^{\beta_i-1}\,dp
 =\Gamma(\beta)D(\beta_1,\ldots,\beta_t).
\]

(3) Let \(\mu\) be the Dirichlet law with density
\(\Gamma(\alpha)\prod_ip_i^{\alpha_i-1}/\prod_i\Gamma(\alpha_i)\) on
\(\Delta\); by (2) with \(\beta_i=\alpha_i\) it is a probability. By
(2) with \(\beta_i=\alpha_i+n_i\),
\[
 \int\prod_ip_i^{n_i}\,d\mu=\frac{\Gamma(\alpha)}{\prod_i\Gamma(\alpha_i)}
 \cdot\frac{\prod_i\Gamma(\alpha_i+n_i)}{\Gamma(\alpha+n)}
 =\frac{\prod_i\alpha_i^{(n_i)}}{\alpha^{(n)}}.
\]
Identify \(p\in\Delta\) with the law on \(\{c_1,\ldots,c_t\}\) giving
\(c_i\) mass \(p_i\). The mixture \(Q(A)=\int p^{\otimes\mathbb N_+}(A)\,\mu(dp)\)
is a probability on sequence space, and by (1) and the last display it
agrees with \(P\) on every cylinder event
\(\{X_1=e_1,\ldots,X_n=e_n\}\). These events together with the empty set
form a \(\pi\)-system generating the product \(\sigma\)-field, so
\(Q=P\). Thus \(\mu\) (as a law on the laws of \(\{c_1,\ldots,c_t\}\),
which form a standard Borel space) is a mixing law for \(P\) in the
sense of \cref{thm:p15-definetti-kernel}, and by
\cref{thm:p15-directing-law-uniqueness} it is the mixing law. The count
formula follows by summing (1) over the
\(n!/(n_1!\cdots n_t!)\) sequences with given counts.
\end{proof}
\end{annexitem}

\begin{annexitem}{annex:p15-w02b-p15-029-isomorphism}{thm:p15-w02b-p15-029-isomorphism}{Invariance of the L-concepts and the distribution lemmas}
\begin{proof}
A state description \(z\) of \(\mathfrak L\) contains, for every atomic
sentence \(A\) that is not an identity, exactly one of \(A\) and
\(\neg A\); in \(\mathfrak L_\infty\) it is a class of such sentences,
in \(\mathfrak L_N\) their conjunction in the fixed order. For an
in-correlation \(C\), let \(C\cdot z\) contain \(C(A)\) exactly when
\(z\) contains \(A\), and \(\neg C(A)\) exactly when \(z\) contains
\(\neg A\). Since \(C\) is one-one and onto, \(A\mapsto C(A)\) is a
bijection of the atomic sentences, so \(C\cdot z\) is a state
description, it is the one constructed from \(z\) by \(C\), and
\(z\mapsto C\cdot z\) is a bijection of the state descriptions with
inverse \(z\mapsto C^{-1}\cdot z\).

(1) We show by induction on the number of connectives and quantifiers
in a sentence \(i\) that \(i\) holds in \(z\) iff \(C(i)\) holds in
\(C\cdot z\). If \(i\) is atomic and not an identity, \(C(i)\) holds in
\(C\cdot z\) iff \(C(i)\in C\cdot z\) iff \(i\in z\). If \(i\) is
\(b=b'\), it holds in every state iff \(b\) and \(b'\) are the same
constant and in none otherwise; since \(C\) is one-one, \(C(b)\) and
\(C(b')\) are the same constant iff \(b\) and \(b'\) are. The steps for
\(\neg\), \(\lor\), \(\land\) are immediate from the truth conditions,
since \(C\) commutes with the connectives. Let \(i\) be \(\forall x\,\varphi\).
By the truth condition for the universal quantifier (every individual
has a constant), \(i\) holds in \(z\) iff \(\varphi[b/x]\) holds in
\(z\) for every constant \(b\). Now \(C(i)\) is \(\forall x\,C(\varphi)\),
and \(C(\varphi)[b/x]\) is \(C(\varphi[C^{-1}(b)/x])\). Each instance
\(\varphi[b/x]\) has fewer quantifiers than \(i\), so by the induction
hypothesis \(C(i)\) holds in \(C\cdot z\) iff \(\varphi[C^{-1}(b)/x]\)
holds in \(z\) for every \(b\), iff \(\varphi[b'/x]\) holds in \(z\) for
every \(b'\) (as \(C^{-1}\) is onto), iff \(i\) holds in \(z\). The
existential quantifier is the dual. Hence
\(R(C(i))=\{C\cdot z:z\in R(i)\}\).

(2) By (1), \(R(i')\) is the image of \(R(i)\) under a bijection of the
set of all state descriptions. A bijection preserves being the whole
set, being empty, inclusion, equality, disjointness, and covering the
whole set by a union. These are L-truth, L-falsity, L-implication,
L-equivalence, L-exclusion and L-disjunctness; factuality is neither
of the first two.

Before (3) we record a fact (\(*\)). By D25-4a, b, every state \(z\)
assigns each constant \(b\) exactly one \(m\) with \(M_mb\) holding in
\(z\); write \(\lambda_z(b)\) for it. An individual distribution holds in
\(z\) iff each of its constants has, under \(\lambda_z\), the kind the
distribution gives it.

(3) Let \(i\) give \(b_l\) the kind \(m_l\) (\(l=1,\ldots,n\)). By
D25-4c, \(M_{m_l}\) is not L-empty, so some state \(z_l\) and constant
\(c_l\) have \(M_{m_l}c_l\) holding in \(z_l\). The truth of \(M_mb\) in
a state depends only on the values there of the atomic sentences
\(P(b)\) for the primitive predicates \(P\) in \(M_m\), all of degree
one. Let \(z\) give \(P(b_l)\) the value that \(z_l\) gives \(P(c_l)\),
and any values to the other atomic sentences; the \(b_l\) are distinct,
so this is consistent, and \(i\) holds in \(z\). Changing the values at
\(b_1\) so that \(b_1\) satisfies some \(M_{m'}\) with \(m'\ne m_1\)
(possible as \(p\geq2\), by the same construction) gives a state in
which \(i\) fails by D25-4b. So \(i\) is factual. If \(i\ne k\) are
individual distributions for the same constants, then, both being in the
fixed order, some \(b\) has different kinds in them, and by D25-4b and
(\(*\)) no state satisfies both. A disjunct L-implies the disjunction,
so \(i\) L-implies \(j\). For the equivalence: if \(k\) is isomorphic
to \(i\), it is a disjunct of \(j\) and so L-implies \(j\). If \(k\)
L-implies \(j\), then \(k\) is a disjunct of \(j\): otherwise \(k\) is
L-exclusive with every disjunct, so \(R(k)\cap R(j)=\varnothing\), and
\(R(k)\ne\varnothing\) gives a state of \(k\) outside \(R(j)\). If \(k\)
is a disjunct of \(j\), it is isomorphic to \(i\), and since isomorphism
is an equivalence relation (with \(C^{-1}\) and composites of
in-correlations), the individual distributions isomorphic to \(k\) are
those isomorphic to \(i\); so \(j\) corresponds to \(k\).

(4) Let \(B\) be the set of the \(n\) constants. If \(k\) differs from
\(C(i)\) only in order, then \(C\) maps \(B\), the constants of \(i\),
onto the constants of \(k\), which form \(B\); so \(C\) restricted to
\(B\) is a permutation \(\sigma\) of \(B\), and \(k\) gives \(\sigma(b)\)
the kind \(i\) gives \(b\). Hence the cardinal numbers agree.
Conversely, if they agree, for each \(m\) the sets of constants of kind
\(m\) in \(i\) and in \(k\) have the same size; a bijection between
them for every \(m\) gives a permutation \(\sigma\) of \(B\), and the
in-correlation equal to \(\sigma\) on \(B\) and the identity elsewhere
maps \(i\) to \(k\) up to order. Individual distributions for \(B\)
correspond one-one to maps \(B\to\{1,\ldots,p\}\), by the fixed order.
Those with fibre sizes \(n_1,\ldots,n_p\) number
\(\binom n{n_1}\binom{n-n_1}{n_2}\cdots\binom{n_p}{n_p}=n!/(n_1!\cdots n_p!)\),
choosing the constants of each kind in turn.

(5) The first equivalence is D27-1 together with the symmetry of
isomorphism. A state description \(\mathfrak Z_k\) has range
\(\{\mathfrak Z_k\}\), and distinct state descriptions are
L-exclusive. So if \(\mathfrak Z_k\) is a disjunct of
\(\mathrm{Str}_i\) it L-implies \(\mathrm{Str}_i\), and otherwise it
L-implies \(\neg\mathrm{Str}_i\) and, being not L-false, does not
L-imply \(\mathrm{Str}_i\). \(\mathrm{Str}_i\) holds in
\(\mathfrak Z_k\) iff \(\mathfrak Z_k\in R(\mathrm{Str}_i)\), which is
the set of disjuncts. Let \(l\) be purely general. Then \(C(l)=l\) for
every \(C\), so by (1) \(R(l)\) is closed under \(z\mapsto C\cdot z\),
and so is \(R(\neg l)\). The range of \(\mathrm{Str}_i\) is the orbit of
\(\mathfrak Z_i\) under these maps, so it lies inside \(R(l)\) or
inside \(R(\neg l)\), according as \(l\) holds in \(\mathfrak Z_i\) or
not. Hence \(R(l)\) is a union of orbits; if it is nonempty, \(l\) is
L-equivalent to the disjunction of the corresponding structure
descriptions.
\end{proof}
\end{annexitem}

\begin{annexitem}{annex:p15-w02b-p15-029-symmetrical}{thm:p15-w02b-p15-029-symmetrical}{Structure descriptions and invariance}
\begin{proof}
(1) By \cref{thm:p15-w02b-p15-029-isomorphism}(5), \(R(\mathrm{Str}_i)\)
is the set of state descriptions isomorphic to \(\mathfrak Z_i\), and
\(m\) of a non-L-false sentence is the sum of \(m\) over its range.

(2) This is D91-1 with D55-3: \(c\) is defined exactly on the pairs with
\(m(e)\ne0\), and for regular \(m\), \(m(e)=0\) iff \(R(e)\) is empty.

(3) \emph{The finite case.} (a) If \(i\) is L-false, so is \(i'\)
(\cref{thm:p15-w02b-p15-029-isomorphism}(2)) and both values are
\(0\). Otherwise \(R(i')=\{C\cdot z:z\in R(i)\}\) by
\cref{thm:p15-w02b-p15-029-isomorphism}(1); \(z\mapsto C\cdot z\) is
one-one and \(C\cdot z\) is isomorphic to \(z\), so
\(m(C\cdot z)=m(z)\) by D90-1, and the sums agree. (b) \(e'\) is
L-false iff \(e\) is. If not,
\(c(h',e')=m(C(e\land h))/m(C(e))=m(e\land h)/m(e)=c(h,e)\) by (a).
(c) Let \(B_h\), \(B_e\) be the sets of constants of \(h\) and \(e\).
The map \(\gamma\) with \(\gamma(b)=C(b)\) for \(b\in B_h\) and
\(\gamma(b)=b\) for \(b\in B_e\) is well defined since
\(B_h\cap B_e=\varnothing\), and one-one since \(C\) is and
\(C(B_h)\), the constants of \(h'\), is disjoint from \(B_e\). Its
domain and range are finite sets of the same size, so their complements
in the set of constants have the same size: equal finite sizes in
\(\mathfrak L_N\), and both are infinite subsets of the constants of
\(\mathfrak L_\infty\). Pairing the two complements in increasing order
of index extends \(\gamma\) to an in-correlation \(C'\); no choice is
involved. Then \(C'(h)=h'\) and \(C'(e)=e\), and (b) for \(C'\) gives
\(c(h',e)=c(h,e)\).

\emph{The infinite case.} Let \(i\) be a sentence of
\(\mathfrak L_\infty\) and \(N_0\) at least the largest index of a
constant in \(i\) or \(C(i)\). For \(N\geq N_0\), the restriction of
\(C\) to the constants of \(i\) is a one-one map between finite subsets
of \(\{a_1,\ldots,a_N\}\); extend it, as in (c), to an in-correlation
\(C_N\) of \(\mathfrak L_N\). Then \(C_N(i)=C(i)\) as sentences, both
occur in \(\mathfrak L_N\), and the finite case gives
\(m_N(C(i))=m_N(i)\) for all \(N\geq N_0\). Two sequences of reals
that agree from some index on either both converge, to the same limit,
or both diverge (T40-21e); this is (a) for \(m_\infty\), including the
clause on existence. For (b), take \(N_0\) for the constants of \(h\)
and \(e\) together and \(C_N\) extending \(C\) on them. For
\(N\geq N_0\), \(c_N(h,e)\) is defined iff \(e\) is not L-false in
\(\mathfrak L_N\), iff \(C_N(e)=e'\) is not L-false in
\(\mathfrak L_N\), iff \(c_N(h',e')\) is defined, and then the two
values are equal by the finite case. So the two sequences of defined
values agree from \(N_0\) on, and the conclusion follows as for (a).
(c) in \(\mathfrak L_\infty\) follows from (b) for the in-correlation
\(C'\) of \(\mathfrak L_\infty\) constructed in (c).

(4) By (1) and D90-1,
\(m(\mathrm{Str}_i)=\sum_{\mathfrak Z\in R(\mathrm{Str}_i)}m(\mathfrak Z)=\zeta_i\,m(\mathfrak Z_i)\).
Since \(\mathfrak Z_i\) L-implies \(\mathrm{Str}_i\),
\(\mathrm{Str}_i\land\mathfrak Z_i\) is L-equivalent to
\(\mathfrak Z_i\), and
\(c(\mathfrak Z_i,\mathrm{Str}_i)=m(\mathfrak Z_i)/(\zeta_i\,m(\mathfrak Z_i))=1/\zeta_i\),
where \(m(\mathfrak Z_i)>0\) by regularity.
\end{proof}
\end{annexitem}

\begin{annexitem}{annex:p15-w02b-p15-029-t92-1}{thm:p15-w02b-p15-029-t92-1}{Individual and statistical distributions as evidence}
\begin{proof}
First let the system be \(\mathfrak L_N\), with all the sentences
occurring in it. Let \(B\) be the set of the \(n\) constants of \(j\),
and \(i_1,\ldots,i_d\) the disjuncts of \(j\), \(i\) among them.

(a) is \cref{thm:p15-w02b-p15-029-isomorphism}(4). (b) \(i'\) is
L-equivalent to \(C(i)\) for some \(C\), and \(m\) is constant on
L-equivalence classes (equal ranges), so (b) is
\cref{thm:p15-w02b-p15-029-symmetrical}(3)(a). (c) The \(i_k\) are
pairwise L-exclusive (\cref{thm:p15-w02b-p15-029-isomorphism}(3)), so
\(R(j)\) is the disjoint union of the \(R(i_k)\) and
\(m(j)=\sum_km(i_k)=d\,m(i)\) by (b). (d) By the proof of
\cref{thm:p15-w02b-p15-029-isomorphism}(4) there is an in-correlation
\(C\) that permutes \(B\), fixes every other constant, and maps \(i\)
to \(i'\) up to order. As \(h\) contains no constant of \(B\),
\(C(h)=h\), so \(C(h\land i)\) is L-equivalent to \(h\land i'\), and
(d) follows from \cref{thm:p15-w02b-p15-029-symmetrical}(3)(a). (e)
\(h\land j\) is L-equivalent to \(\bigvee_k(h\land i_k)\), whose
disjuncts are pairwise L-exclusive, and each has the value
\(m(h\land i)\) by (d). (f) \(c(i,e)=m(e\land i)/m(e)=m(e\land i')/m(e)=c(i',e)\)
by (d) with \(e\) for \(h\). (g) \(c(j,e)=m(e\land j)/m(e)=d\,m(e\land i)/m(e)\)
by (e). (h) \(i\) and \(i'\) are factual
(\cref{thm:p15-w02b-p15-029-isomorphism}(3)), so \(m(i)=m(i')>0\) by
regularity and (b); with (d),
\(c(h,i)=m(i\land h)/m(i)=m(i'\land h)/m(i')=c(h,i')\). (i)
\(c(h,j)=m(j\land h)/m(j)=d\,m(i\land h)/(d\,m(i))=c(h,i)\) by (c),
(e). (j) is (c) with (a), and (k) is (g) with (a).

For \(\mathfrak L_\infty\), let \(N_0\) be the largest index of a
constant in the sentences of the item. For \(N\geq N_0\) the identity
holds in \(\mathfrak L_N\) with a constant factor (\(1\), \(d\), or the
multinomial coefficient), and the two sides are defined for the same
\(N\): an \(m\)-value always, \(c_N(\cdot,e)\) when \(e\) is not
L-false in \(\mathfrak L_N\), and \(c_N(h,i)\), \(c_N(h,i')\),
\(c_N(h,j)\) always, since \(i\), \(i'\), \(j\) are factual. If
\(x_N=\lambda y_N\) for all \(N\geq N_0\) with a fixed \(\lambda>0\),
then \((x_N)\) converges iff \((y_N)\) does, and
\(\lim x_N=\lambda\lim y_N\) (T40-21c, e). For a fitting sequence and
nongeneral \(h,e\), the \(m\)-values and \(c\)-values are constant from
the first system containing the sentences on
(\cref{thm:p15-w02b-p15-029-null-fitting}(2), (3)), so all the limits
exist.
\end{proof}
\end{annexitem}

\begin{annexitem}{annex:p15-w02b-p15-029-t92-3}{thm:p15-w02b-p15-029-t92-3}{A statistical distribution and its part}
\begin{proof}
Work in \(\mathfrak L_N\) with \(N\) at least the largest index of the
\(s+s'\) constants. Let \(B\) and \(B'\) be the constants of \(i\) and
\(i'\), and use \(\lambda_z\) from (\(*\)) in the proof of
\cref{thm:p15-w02b-p15-029-isomorphism}. By that lemma (3) and (4), a
statistical distribution for a set \(A\) with cardinal numbers
\(n_1,\ldots,n_p\) holds in \(z\) iff
\(\lvert\{b\in A:\lambda_z(b)=m\}\rvert=n_m\) for every \(m\), since its
disjuncts are exactly the individual distributions for \(A\) with those
numbers. The conjunction \(i\land i'\) differs only in order from an
individual distribution \(k\) for \(B\cup B'\) with cardinal numbers
\(s_m+s'_m\), and \(e\) is the statistical distribution corresponding
to \(k\).

(a) is \cref{thm:p15-w02b-p15-029-isomorphism}(4) applied to \(j\),
\(j'\), \(e\). (b) is \cref{thm:p15-w02b-p15-029-t92-1}(j) for \(k\)
and \(e\), with \(m(k)=m(i\land i')\).

(c) If \(j\land j'\) holds in \(z\), \(B\) has \(s_m\) and \(B'\) has
\(s'_m\) constants of kind \(m\), and as \(B\cap B'=\varnothing\),
\(B\cup B'\) has \(s_m+s'_m\); so \(e\) holds. (d) If \(e\land j\)
holds, \(B\cup B'\) has \(s_m+s'_m\) and \(B\) has \(s_m\) constants of
kind \(m\), so the disjoint \(B'\) has \(s'_m\); so \(j'\) holds. (e)
From (c) and (d). (f) \(i\) L-implies \(j\); so if \(e\land i\) holds,
\(j\) holds, hence \(j'\) by (d), and \(i\land j'\) holds. If
\(i\land j'\) holds, \(j\land j'\) holds, hence \(e\) by (c).

(g) The first equation is (f). Since \(i\) contains no constant of
\(j'\), \cref{thm:p15-w02b-p15-029-t92-1}(e) with \(h=i\), the
distribution \(i'\) and the statistical distribution \(j'\) gives
\(m(i\land j')=(s'!/\prod_ms'_m!)\,m(i\land i')\).

(h) Since \(j'\) contains no constant of \(j\),
\cref{thm:p15-w02b-p15-029-t92-1}(e) with \(h=j'\), the distribution
\(i\) and \(j\) gives the first equation; the second is (g).

(i) \(i\land i'\) is L-equivalent to the individual distribution \(k\),
which is factual, so \(m(i\land i')>0\). By (g) and (b),
\[
 c(i,e)=\frac{m(e\land i)}{m(e)}
 =\frac{s'!}{\prod_ms'_m!}\cdot\frac{\prod_m(s_m+s'_m)!}{(s+s')!}.
\]
Since \(\genfrac{[}{]}{0pt}{}{s_m+s'_m}{s_m}=(s_m+s'_m)!/s'_m!\) and
\(\genfrac{[}{]}{0pt}{}{s+s'}{s}=(s+s')!/s'!\), the quotient of the
product of the first by the second is the same number.

(j) By (e), \(m(e\land j)=m(j\land j')\); by (h) and (b),
\[
 c(j,e)=\frac{s!}{\prod_ms_m!}\cdot\frac{s'!}{\prod_ms'_m!}\cdot
 \frac{\prod_m(s_m+s'_m)!}{(s+s')!},
\]
which is \((s!/\prod_ms_m!)\,c(i,e)\) by (i), and equals
\(\prod_m\binom{s_m+s'_m}{s_m}/\binom{s+s'}{s}\) on writing each
binomial coefficient as a quotient of factorials.

The values in (i) and (j) do not depend on \(N\) or on the symmetrical
function. In \(\mathfrak L_\infty\), \(e\) is nongeneral and
factual, so \(c_N(i,e)\) and \(c_N(j,e)\) are defined for every \(N\)
in which the constants occur, and the constant sequences converge to
these numbers (T40-21d).
\end{proof}
\end{annexitem}
